\chapter{False Tate extensions}

Let $p$ be prime, let $K/\Qp$ be finite, and fix an algebraic closure
$\overline K$. Let $\pi$ be a uniformizer of $K$. All field inclusions below
are taken inside this fixed closure. Valuations of finite local fields are
normalized to have value group $\mathbb Z$.

\begin{definition}[Compatible radicals and roots of unity]
\label{ft:compatible-roots}
\lean{FalseTateExtension.exists_compatible_radicals, FalseTateExtension.exists_compatible_primitive_roots, FalseTateExtension.exists_compatible_roots, FalseTateExtension.radical, FalseTateExtension.radical_zero, FalseTateExtension.radical_succ_pow, FalseTateExtension.radical_pow, FalseTateExtension.rootOfUnity, FalseTateExtension.rootOfUnity_zero, FalseTateExtension.rootOfUnity_succ_pow, FalseTateExtension.isPrimitiveRoot_rootOfUnity, FalseTateExtension.rootOfUnity_pow}
\leanok
\source{liu-2008-lattices:page-0022}
\dcref{Liu, Section 5.1}
Choose $\pi_0=\pi$ and $\pi_{m+1}^p=\pi_m$ for $m\geq0$.
Choose primitive $p^n$-th roots of unity $\zeta_n$ for $n\geq1$ such that
$\zeta_{n+1}^p=\zeta_n$. Put $\zeta_0=1$.
\end{definition}
\begin{proof}
Such systems exist in every algebraically closed field of characteristic
zero. Starting at $\pi_0=\pi$, recursively choose a root of
$X^p-\pi_n$. Induction gives $\pi_n^{p^n}=\pi$, including $\pi=0$.
Choose a primitive $p$-th root $\zeta_1$: the polynomial $X^p-1$ has
$p$ distinct roots, and any root other than $1$ has order $p$.
Having chosen $\zeta_n$ for $n\geq1$, choose a root $\eta$ of
$X^p-\zeta_n$ and set $\zeta_{n+1}=\eta$.
Then $\eta^{p^{n+1}}=1$, so its order divides $p^{n+1}$ and hence
is $p^r$ for some $r\leq n+1$. If $r\leq n$, then
$\zeta_n^{p^{n-1}}=\eta^{p^n}=1$, contradicting primitivity of
$\zeta_n$. Thus $r=n+1$. This proves primitivity at every level;
the step from $\zeta_0=1$ is the initial choice. The argument includes
$p=2$.
\end{proof}

\begin{definition}[Finite layers and infinite extensions]
\label{ft:extensions}\uses{ft:compatible-roots}
\lean{FalseTateExtension.radicalLayer, FalseTateExtension.radical_mem_radicalLayer, FalseTateExtension.radicalLayer_zero, FalseTateExtension.radicalLayer_le_succ, FalseTateExtension.radicalLayer_mono, FalseTateExtension.radicalLayer_finiteDimensional, FalseTateExtension.kummerExtension, FalseTateExtension.radicalLayer_le_kummerExtension, FalseTateExtension.kummerExtension_eq_iSup_succ, FalseTateExtension.mem_kummerExtension_iff, FalseTateExtension.cyclotomicLayer, FalseTateExtension.rootOfUnity_mem_cyclotomicLayer, FalseTateExtension.cyclotomicLayer_zero, FalseTateExtension.cyclotomicLayer_le_succ, FalseTateExtension.cyclotomicLayer_mono, FalseTateExtension.cyclotomicLayer_finiteDimensional, FalseTateExtension.cyclotomicExtension, FalseTateExtension.cyclotomicLayer_le_cyclotomicExtension, FalseTateExtension.cyclotomicExtension_eq_iSup_succ, FalseTateExtension.mem_cyclotomicExtension_iff, FalseTateExtension.falseTateLayer, FalseTateExtension.falseTateLayer_zero_zero, FalseTateExtension.radical_mem_falseTateLayer, FalseTateExtension.rootOfUnity_mem_falseTateLayer, FalseTateExtension.falseTateLayer_mono, FalseTateExtension.falseTateLayer_directed, FalseTateExtension.falseTateLayer_le_diagonal, FalseTateExtension.falseTateLayer_finiteDimensional, FalseTateExtension.falseTateExtension, FalseTateExtension.falseTateLayer_le_falseTateExtension, FalseTateExtension.falseTateExtension_eq_iSup_diagonal, FalseTateExtension.mem_falseTateExtension_iff_diagonal, FalseTateExtension.mem_falseTateExtension_iff, FalseTateExtension.falseTateExtension_eq_iSup_succ, FalseTateExtension.mem_falseTateExtension_iff_pos, FalseTateExtension.falseTateExtension_eq_sup}
\leanok
\source{liu-2008-lattices:page-0022}
\dcref{Liu, Section 5.1}
For $m,n\geq1$, set
\[
 K^{\mathrm{FT}}_{m,n}=K(\pi_m,\zeta_n),\qquad
 K^{\mathrm{FT}}_\infty=\bigcup_{m,n\geq1}K^{\mathrm{FT}}_{m,n}.
\]
Set $K^{\mathrm{Kum}}=\bigcup_mK(\pi_m)$ and
$K^{\mathrm{cyc}}=\bigcup_nK(\zeta_n)$.
The unions are directed by the two indices, and the false Tate extension
is the compositum $K^{\mathrm{Kum}}K^{\mathrm{cyc}}$.
\end{definition}
\begin{proof}
The successive power relations put $\pi_m$ in $K(\pi_{m+1})$ and
$\zeta_n$ in $K(\zeta_{n+1})$. Repeated inclusion proves monotonicity
in each index. The generators satisfy the monic equations
$X^{p^m}-\pi=0$ and $X^{p^n}-1=0$, respectively. Each simple layer
is therefore finite over $K$, and adjoining both generators gives a
finite layer as a tower of two finite simple extensions. In
characteristic zero these extensions are separable.

For two pairs of indices take the coordinatewise maximum; the resulting
layer contains both original layers. Hence their union is a field:
each field operation on finitely many elements can be performed in one
common layer. The diagonal layers are cofinal, because the layer at
$(m,n)$ is contained in the layer at $(r,r)$ for $r=\max(m,n)$.
Allowing index zero changes none of these infinite fields, by the same
cofinality with indices at least one. Each finite two-generator layer
lies in the compositum of the two infinite simple towers. Conversely,
the union of the two-generator layers contains every generator of both
towers, and thus their compositum. These two inclusions prove the final
identity.
\end{proof}

\begin{lemma}[Galois finite layers and their directed unions]
\label{ft:galois-layers}
\source{liu-2008-lattices:page-0022}
\dcref{Liu, Section 5.1; finite splitting-field argument at every prime}
\uses{ft:extensions}
\lean{FalseTateExtension.cyclotomicLayer_isGalois, FalseTateExtension.radical_ne_zero, FalseTateExtension.adjoin_rootSet_eq_falseTateLayer, FalseTateExtension.falseTateLayer_isSplittingField, FalseTateExtension.falseTateLayer_isGalois, FalseTateExtension.cyclotomicExtension_isGalois, FalseTateExtension.falseTateExtension_isGalois}
\leanok
For every prime $p$ and every characteristic-zero field $K$, the
cyclotomic layers $K(\zeta_n)/K$ are finite Galois. If $\pi\ne0$,
the diagonal fields $K(\pi_n,\zeta_n)/K$ are finite Galois. Their
directed infinite cyclotomic and false Tate unions are Galois over $K$.
No degree or ramification formula is asserted here.
\end{lemma}
\begin{proof}
Put $d=p^n$. The roots of $X^d-1$ are exactly the powers of the
primitive root $\zeta_n$. Thus $K(\zeta_n)$ is its splitting field.
For the radical polynomial put $\alpha=\pi_n$; then $\alpha^d=\pi$,
so $\alpha\ne0$. Its roots are precisely $\alpha\zeta_n^j$:
division of any root by $\alpha$ gives a $d$-th root of unity, and
conversely every displayed element is a root. All these roots lie in
$K(\alpha,\zeta_n)$. Conversely the field generated by the roots
contains $\alpha$ and $\alpha\zeta_n$, and hence their ratio
$\zeta_n$. Therefore the diagonal field is exactly the splitting
field of $X^d-\pi$. This includes $n=0$. In characteristic zero both
splitting-field extensions are separable, hence finite Galois.

Every element of either directed union lies in a finite layer. Its
minimal polynomial over $K$ splits in that layer by normality, and
therefore splits in the union. The unions are thus normal algebraic
extensions. Characteristic zero supplies separability, so they are
Galois. The diagonal cofinality in the preceding definition identifies
the second union with the actual two-index false Tate extension.
\end{proof}

\begin{lemma}[Finite cyclotomic and Kummer coordinates]
\label{ft:finite-coordinates}
\source{liu-2008-lattices:page-0022}
\dcref{Liu, Section 5.1; finite coordinate calculation}
\uses{ft:galois-layers,ft:compatible-roots}
\lean{FalseTateExtension.layerRadical, FalseTateExtension.layerRootOfUnity, FalseTateExtension.coe_layerRadical, FalseTateExtension.coe_layerRootOfUnity, FalseTateExtension.isPrimitiveRoot_layerRootOfUnity, FalseTateExtension.layerRadical_pow, FalseTateExtension.layerRadical_ne_zero, FalseTateExtension.finiteCyclotomicCharacter, FalseTateExtension.finiteCyclotomicCharacter_spec, FalseTateExtension.layerRootUnit, FalseTateExtension.coe_layerRootUnit, FalseTateExtension.isPrimitiveRoot_layerRootUnit, FalseTateExtension.kummerRatio, FalseTateExtension.coe_kummerRatio, FalseTateExtension.kummerRatio_mem_zpowers, FalseTateExtension.finiteKummerExponent, FalseTateExtension.finiteKummerExponent_spec, FalseTateExtension.layerRootOfUnity_pow_injective, FalseTateExtension.finiteKummerExponent_unique, FalseTateExtension.finiteCyclotomicCharacter_unique, FalseTateExtension.finiteKummerExponent_mul, FalseTateExtension.finiteLayer_aut_ext, FalseTateExtension.finiteCoordinates_injective}
\leanok
Let $\pi\ne0$, $n\geq0$, $q=p^n$, and $E=K(\pi_n,\zeta_n)$.
There are actual maps
\[
 \chi_n:\Gal(E/K)\longrightarrow(\mathbb Z/q\mathbb Z)^\times,
 \qquad b_n:\Gal(E/K)\longrightarrow\mathbb Z/q\mathbb Z
\]
with $\chi_n$ a group homomorphism, characterized by
\[
 g(\zeta_n)=\zeta_n^{\chi_n(g)},\qquad
 g(\pi_n)=\zeta_n^{b_n(g)}\pi_n.
\]
They satisfy
\[
 b_n(gh)=b_n(g)+\chi_n(g)b_n(h),
\]
and the map $g\mapsto(b_n(g),\chi_n(g))$ is injective. This includes
$p=2$ and $n=0$.
\end{lemma}
\begin{proof}
An automorphism preserves primitive $q$-th roots, so it sends $\zeta_n$
to a unique power with exponent a unit modulo $q$. Evaluating a product
of automorphisms on $\zeta_n$ and using exponent uniqueness proves that
$\chi_n$ is a homomorphism. Since $\pi_n\ne0$ and both $g(\pi_n)$ and
$\pi_n$ have $q$-th power $\pi$, their ratio is a $q$-th root of unity.
It is a unique power of $\zeta_n$ with exponent modulo $q$, defining
$b_n(g)$.

Apply $g$ to $h(\pi_n)=\zeta_n^{b_n(h)}\pi_n$ and substitute both
evaluation formulas. This gives
\[
 (gh)(\pi_n)=\zeta_n^{\chi_n(g)b_n(h)+b_n(g)}\pi_n.
\]
Cancel the nonzero factor $\pi_n$ and use the order $q$ of $\zeta_n$
to obtain the cocycle identity modulo $q$. Automorphisms with the same
coordinates agree on $\pi_n$ and $\zeta_n$, so they agree on the field
these elements generate. This proves injectivity. When $n=0$, the root
of unity is $1$, the radical lies in $K$, and both coordinate groups are
trivial, consistent with the same argument.
\end{proof}

\begin{lemma}[Finite cyclotomic groups at odd primes]
\label{ft:odd-finite-cyclotomic-cyclic}
\source{liu-2008-lattices:page-0023}
\dcref{Finite cyclic form of the cyclotomic group argument in Lemma 5.1.2}
\uses{ft:extensions,ft:galois-layers,ft:compatible-roots}
\lean{FalseTateExtension.cyclotomicLayer_isCyclic}
\leanok
Assume $p\geq3$. For every $m\geq0$, the actual group
$\operatorname{Gal}(K(\zeta_m)/K)$ is finite cyclic.
\end{lemma}
\begin{proof}
The actual cyclotomic character embeds this group in
$(\mathbb Z/p^m\mathbb Z)^\times$: its kernel fixes the generator
$\zeta_m$ and hence the whole field. The unit group modulo an odd
prime power is cyclic, including the trivial modulus $p^0=1$.
A subgroup of a cyclic group is cyclic. The source's infinite
principal-unit argument is thus replaced here by its finite cyclic
consequence, available for every actual finite cyclotomic field.
\end{proof}

\begin{lemma}[The first radical is not cyclotomic at an odd prime]
\label{ft:first-radical-not-cyclotomic}
\source{liu-2008-lattices:page-0022,liu-2008-lattices:page-0023}
\dcref{Liu, Lemma 5.1.2, first nontrivial radical step; finite cyclic adaptation}
\uses{ft:odd-finite-cyclotomic-cyclic,ft:radical-layer-ramification,ft:compatible-roots,ft:extensions}
\lean{FalseTateExtension.radical_one_not_mem_cyclotomicExtension}
\leanok
Assume $p\geq3$. Then $\pi_1\notin K^{\mathrm{cyc}}$.
\end{lemma}
\begin{proof}
Suppose otherwise and choose $m$ with
$F=K(\pi_1)\subseteq C_m=K(\zeta_m)$.
Eisenstein irreducibility gives $[F:K]=p$.
The cyclic extension $C_m/K$ is abelian, so $F/K$ is Galois.
The polynomial $X^p-\pi$ therefore splits in $F$; ratios of its
roots give a primitive $p$th root of unity in $F$.
The degree $[K(\zeta_1):K]$ divides $p$ and is at most $p-1$,
so it equals one. Thus $\zeta_1\in K$.

Choose the least $r\leq m$ such that $C_r\ne K$.
It exists since $F\subseteq C_m$ and $[F:K]=p$.
Then $r\geq2$ and $\zeta_{r-1}\in K$.
The actual group of $C_r/K$ embeds in the $p$th roots of unity
by its action on $\zeta_r$, since $\zeta_r^p\in K$.
Its nontrivial degree therefore divides $p$ and equals $p$.
A finite cyclic group has a unique subgroup of index $p$:
if a generator is $g$, that subgroup is generated by $g^p$.
Galois correspondence inside $C_m$ now gives $F=C_r$.

Choose $1\ne\sigma\in\operatorname{Gal}(F/K)$.
There are $b,c\in(\mathbb Z/p\mathbb Z)^\times$ with
\[
 \sigma(\pi_1)=\zeta_1^b\pi_1,\qquad
 \sigma(\zeta_r)=\zeta_1^c\zeta_r.
\]
Neither exponent vanishes because each displayed element generates
$F$ over $K$. Write
$\pi_1=\sum_{i=0}^{p-1}a_i\zeta_r^i$ in the power basis of $C_r/K$.
Since $\zeta_1\in K$, applying $\sigma$ and comparing coefficients
gives $a_i(\zeta_1^{ci}-\zeta_1^b)=0$ for each $i$.
There is exactly one residue $i$ with $ci=b$, and it is nonzero.
As $\pi_1\ne0$, its coefficient is nonzero and
$\pi_1=a\zeta_r^i$ for some $a\in K^\times$.
Taking $p$th powers yields
$\pi=a^p\zeta_{r-1}^i$. A root of unity is a unit, so normalized
valuation on the original field $K$ gives $1=pv_K(a)$, impossible
for the integer $v_K(a)$.
This reverse power-basis comparison gives the same unit-versus-radical
contradiction as Liu's first step, using finite cyclicity in place
of an infinite logarithm classification.
\end{proof}

\begin{lemma}[The Galois structure]
\label{ft:galois-structure}\uses{ft:first-radical-not-cyclotomic,ft:affine-embedding,ft:affine-additive-kernel,ft:padic-closed-additive-subgroups,ft:finite-coordinates,fw:profinite-galois-foundations}\lean{FalseTateExtension.kummerExtension_inf_cyclotomicExtension, FalseTateExtension.falseTateGaloisEquiv, FalseTateExtension.cyclotomicRelativeGaloisEquiv, FalseTateExtension.kummerRelativeGaloisEquiv, FalseTateExtension.falseTateGaloisEquiv_left, FalseTateExtension.falseTateGaloisEquiv_right, FalseTateExtension.falseTateGaloisEquiv_mem_translations, FalseTateExtension.falseTateGaloisEquiv_mem_units}
\leanok
\source{liu-2008-lattices:page-0022,liu-2008-lattices:page-0023}
\dcref{Liu, Lemma 5.1.2}
Assume $p\geq3$. Then
\[
 K^{\mathrm{Kum}}\cap K^{\mathrm{cyc}}=K.
\]
Write $H_K=\Gal(K^{\mathrm{cyc}}/K)$, viewed through the cyclotomic
character as a closed subgroup of $\Zp^\times$. Restriction and the action
on the compatible radicals give topological isomorphisms
\[
 \Gal(K^{\mathrm{FT}}_\infty/K^{\mathrm{cyc}})\simeq\Zp(1),\qquad
 \Gal(K^{\mathrm{FT}}_\infty/K^{\mathrm{Kum}})\simeq H_K,
\]
and
\[
 \Gal(K^{\mathrm{FT}}_\infty/K)\simeq\Zp(1)\rtimes H_K,
\]
where $H_K$ acts on $\Zp(1)$ through the cyclotomic character.
\end{lemma}
\begin{proof}
Use the actual closed affine realization $G$ of the original
Galois group, with translation kernel $M\subseteq\mathbb Z_p$
and actual unit projection $H_K$.
If $M\subseteq p\mathbb Z_p$, every automorphism of the false Tate
field fixing its cyclotomic subfield has translation coordinate
zero modulo $p$, so fixes $\pi_1$ by the finite coordinate formula.
Its fixed field is the actual cyclotomic subfield by infinite
Galois correspondence. This would put $\pi_1$ in that subfield,
contrary to the preceding lemma.
Thus $M\nsubseteq p\mathbb Z_p$.
The closed-additive-subgroup classification gives
$M=p^a\mathbb Z_p$ or $M=0$; the noncontainment forces $a=0$.
Hence $M=\mathbb Z_p$.

For each $h\in H_K$, choose $(b,h)\in G$.
The element $(-b,1)$ belongs to the now full translation kernel,
and $(-b,1)(b,h)=(0,h)$.
Thus the actual pure-unit subgroup is exactly $\{0\}\rtimes H_K$,
and every pair $(b,h)$ with $b\in\mathbb Z_p$, $h\in H_K$
belongs to $G$. Restricting the original affine homeomorphism
therefore identifies $G$ with $\mathbb Z_p\rtimes H_K$,
continuously in both directions. Its multiplication law is the
semidirect law with the cyclotomic action.
The original subgroups fixing the cyclotomic and radical fields
are respectively the full translations and the pure units.
They generate $G$, so the intersection of their actual fixed
fields is the fixed field of $G$, namely $K$.
Restriction of the same homeomorphism to the two actual fixing
subgroups gives both asserted topological group isomorphisms.
\end{proof}


\begin{definition}[The Kummer generator]
\label{ft:tau}\uses{ft:galois-structure}\lean{FalseTateExtension.relativeTau, FalseTateExtension.relativeTau_extensionRadical, FalseTateExtension.cyclotomicRelativeGaloisEquiv_relativeTau, FalseTateExtension.relativeTau_unique, FalseTateExtension.denseRange_relativeTau_zpow}
\leanok
\source{liu-2008-lattices:page-0023}
\dcref{Liu, discussion after Lemma 5.1.2}
Assume $p\geq3$. The element $\tau\in\Gal(K^{\mathrm{FT}}_\infty/K^{\mathrm{cyc}})$
is specified by
\[
 \tau(\pi_m)=\zeta_m\pi_m\quad(m\geq1).
\]
It corresponds to $1\in\Zp$ under the isomorphism determined by the
chosen roots of unity, and is a topological generator.
\end{definition}
\begin{proof}
Under the preceding isomorphism choose the actual element
$(1,1)\in\mathbb Z_p\rtimes H_K$. It has trivial cyclotomic
character, so fixes the actual cyclotomic field.
The finite Kummer coordinate formula gives its value on every
$\pi_m$ as $\zeta_m\pi_m$.
These values, together with fixing all roots of unity, determine
its action on every finite false Tate layer and hence uniquely
on their union. Its integral powers correspond to the integers
inside $\mathbb Z_p$, whose image is dense. It is therefore a
topological generator of the actual translation subgroup.
\end{proof}


\section{Finite residue degree at every prime}

The argument in this section applies to every prime, including $2$.
It uses the action on radicals rather than a splitting of the Galois
extension into Kummer and cyclotomic factors.

\begin{lemma}[The p-adic cyclotomic polynomials]
\label{ft:padic-cyclotomic-polynomial}
\source{serre-1979-local-fields:page-0085}
\dcref{Serre, Chapter IV, Section 4, proof of Proposition 17}
\lean{PadicInt.coeff_zero_cyclotomic_prime_pow_comp_X_add_one, PadicInt.monic_cyclotomic_prime_pow_comp_X_add_one, PadicInt.natDegree_cyclotomic_prime_pow_comp_X_add_one, PadicInt.p_not_mem_span_p_sq, PadicInt.cyclotomic_prime_pow_comp_X_add_one_isEisensteinAt, PadicInt.irreducible_cyclotomic_prime_pow_comp_X_add_one, PadicInt.cyclotomic_prime_pow_comp_X_add_one_isEisensteinAt_of_pos, Padic.irreducible_cyclotomic_prime_pow_comp_X_add_one, Padic.irreducible_cyclotomic_prime_pow, Padic.minpoly_primitive_root, Padic.finrank_adjoin_primitive_root, FalseTateExtension.finrank_cyclotomicLayer_padic, FalseTateExtension.finrank_cyclotomicLayer_padic_zero}
\leanok
For every prime $p$ and integer $m\geq1$, the polynomial
$\Phi_{p^m}(1+T)\in\Zp[T]$ is Eisenstein at $(p)$, with constant
coefficient $p$ and degree $(p-1)p^{m-1}$.
For every $m\geq0$, $\Phi_{p^m}$ is irreducible over $\Qp$.
Consequently a primitive $p^m$-th root $\zeta_m$ generates an extension
of $\Qp$ of degree $\varphi(p^m)$, including degree one when $m=0$.
\end{lemma}
\begin{proof}
For $m\geq1$ the prime-power cyclotomic identity gives
\[
 \Phi_{p^m}(1+T)=\sum_{j=0}^{p-1}(1+T)^{jp^{m-1}}.
\]
This polynomial is monic of the stated degree and has constant term
$p$. In characteristic $p$, put $X=T^{p^{m-1}}$. The reduction of
the right side is $\sum_{j=0}^{p-1}(1+X)^j=X^{p-1}$:
multiplying by $X$ gives $(1+X)^p-1=X^p$, and cancellation by
$X$ in $\mathbb F_p[X]$ proves the identity. Thus every coefficient
below the leading term is divisible by $p$, whereas the constant
term $p$ is not divisible by $p^2$. This is the Eisenstein criterion
over the discrete valuation ring $\Zp$. Gauss's lemma proves
irreducibility over its fraction field $\Qp$. Translation by one
is an invertible polynomial substitution, so $\Phi_{p^m}$ is
irreducible as well. For $m=0$ it is the linear polynomial $T-1$.
Finally $\zeta_m$ is a root of this monic irreducible polynomial,
so it is its minimal polynomial, and its degree is the extension
degree. The prime-power degree formula gives $\varphi(p^m)$.
\end{proof}

\begin{lemma}[Total ramification of the rational cyclotomic layers]
\label{ft:cyclotomic-total-ramification}
\source{serre-1979-local-fields:page-0085}
\dcref{Serre, Chapter IV, Section 4, Proposition 17}
\uses{ft:padic-cyclotomic-polynomial,fw:norm-valuation,fw:local-dvr-degree-identity}
\lean{FalseTateExtension.padic_valuationSubring_toSubring, FalseTateExtension.padicUnitBallEquiv, FalseTateExtension.padicUnitBall_isDiscreteValuationRing, FalseTateExtension.padicUnitBall_normalizedValuation_p, FalseTateExtension.cyclotomicLayerRoot, FalseTateExtension.isPrimitiveRoot_cyclotomicLayerRoot, FalseTateExtension.cyclotomicLayer_padic_isCyclotomicExtension, FalseTateExtension.norm_cyclotomicLayerRoot_sub_one, FalseTateExtension.cyclotomicValuationSubring_isDiscreteValuationRing, FalseTateExtension.cyclotomicLayer_normalizedValuation_pow, FalseTateExtension.cyclotomicLayer_residueDegree, FalseTateExtension.cyclotomicLayer_normalizedValuation_sub_one, FalseTateExtension.cyclotomicLayerRoot_sub_one_mem, FalseTateExtension.cyclotomicUniformizer, FalseTateExtension.irreducible_cyclotomicUniformizer, FalseTateExtension.span_cyclotomicUniformizer, FalseTateExtension.cyclotomicLayer_ramificationIdx, FalseTateExtension.cyclotomicLayer_ramificationIdx_eq_totient, FalseTateExtension.cyclotomicLayer_spectral_residueDegree, FalseTateExtension.irreducible_cyclotomicUniformizer_spectral, FalseTateExtension.cyclotomicLayer_spectral_ramificationIdx}
\leanok
For every prime $p$ and $m\geq1$, put $E_m=\Qp(\zeta_m)$ for a
primitive $p^m$-th root of unity. Then $E_m/\Qp$ is totally ramified,
and $\zeta_m-1$ is a uniformizer of its actual valuation ring.
\end{lemma}
\begin{proof}
Put $\lambda=\zeta_m-1$ and $d=\varphi(p^m)$. Its monic minimal
polynomial is $\Phi_{p^m}(1+T)$, by irreducibility and translation.
The constant coefficient is $p$, so
$N_{E_m/\Qp}(\lambda)=(-1)^d p$ and
$v_{\Qp}(N(\lambda))=1$. The element $\lambda$ is integral, being
the difference of the integral root of unity and one.
If $f$ is the residue degree, the normalized norm identity gives
$1=fv_{E_m}(\lambda)$. Here $f$ is a positive integer and
$v_{E_m}(\lambda)$ is a nonnegative integer. Their product is one,
so both are one. Thus $\lambda$ is a uniformizer, and the
fundamental degree identity gives $e(E_m/\Qp)=d$, proving total
ramification. This proves the ramification assertion through the
norm of the Eisenstein generator.
\end{proof}

\begin{lemma}[Finite rational cyclotomic characters]
\label{ft:finite-cyclotomic-character}
\uses{ft:padic-cyclotomic-polynomial,ft:galois-layers}
\lean{FalseTateExtension.cyclotomicLayerCharacter, FalseTateExtension.cyclotomicLayerCharacter_spec, FalseTateExtension.cyclotomicLayerCharacter_unique, FalseTateExtension.cyclotomicLayerCharacter_injective, FalseTateExtension.card_cyclotomicLayer_galois_eq_finrank, FalseTateExtension.card_cyclotomicLayer_galois, FalseTateExtension.card_cyclotomicLayer_character_units, FalseTateExtension.cyclotomicLayerCharacter_bijective, FalseTateExtension.cyclotomicLayerCharacter_surjective, FalseTateExtension.cyclotomicLayerGaloisEquivUnits, FalseTateExtension.cyclotomicLayerGaloisEquivUnits_apply, FalseTateExtension.cyclotomicLayerGaloisEquivUnits_spec, FalseTateExtension.cyclotomicLayerGaloisEquivUnits_symm_spec}
\leanok
For every prime $p$ and $m\geq0$, the actual cyclotomic character
of $E_m=\Qp(\zeta_m)$ is an isomorphism
\[
 \Gal(E_m/\Qp)\simeq(\mathbb Z/p^m\mathbb Z)^\times.
\]
Here the character sends $\sigma$ to the unique unit $u$ satisfying
$\sigma(\zeta_m)=\zeta_m^u$.
\end{lemma}
\begin{proof}
The root $\zeta_m$ generates $E_m$, so an automorphism whose
character is one fixes a field generator and is the identity.
Thus the character is injective. The extension is finite Galois,
and its degree is $\varphi(p^m)$ by the polynomial lemma.
Its automorphism group and the indicated unit group consequently
both have $\varphi(p^m)$ elements. An injective map between these
finite sets of the same cardinality is surjective. The character
is multiplicative by its defining action on the root, so it is
a group isomorphism. At $m=0$, both groups are trivial, and the
same degree and cardinality argument applies.
\end{proof}

\begin{lemma}[The infinite rational cyclotomic character]
\label{ft:infinite-cyclotomic-character}
\uses{ft:finite-cyclotomic-character,ft:padic-cyclotomic-polynomial,ft:cyclotomic-projection-image}
\lean{FalseTateExtension.cyclotomicLayerRestriction, FalseTateExtension.cyclotomicLayerRestriction_commutes, FalseTateExtension.cyclotomicLayerRestriction_surjective, FalseTateExtension.continuous_cyclotomicLayerRestriction, FalseTateExtension.cyclotomicLayerCharacter_restrict, FalseTateExtension.cyclotomicCharacterFiber, FalseTateExtension.isClosed_cyclotomicCharacterFiber, FalseTateExtension.cyclotomicCharacterFiber_nonempty, FalseTateExtension.cyclotomicCharacterFiber_succ_subset, FalseTateExtension.cyclotomicExtensionCharacter_surjective_padic, FalseTateExtension.cyclotomicExtensionGaloisEquivUnits, FalseTateExtension.cyclotomicExtensionGaloisEquivUnits_apply, FalseTateExtension.cyclotomicLayer_padic_degrees_unbounded, FalseTateExtension.not_finiteDimensional_cyclotomicExtension_padic}
\leanok
For every prime $p$, the actual cyclotomic character of
$C=\Qp(\mu_{p^\infty})$ is a continuous group isomorphism
\[
 \chi:\Gal(C/\Qp)\simeq\Zp^\times,
\]
with continuous inverse. The extension $C/\Qp$ has infinite degree.
\end{lemma}
\begin{proof}
The actual character is a continuous closed embedding, since the
chosen roots generate $C$. Fix $u\in\Zp^\times$. For each $m\geq0$,
let $D_m$ be the subset of $\Gal(C/\Qp)$ on which reduction of the
character modulo $p^m$ equals the reduction of $u$.
The finite character is surjective, and every automorphism of
$E_m=\Qp(\zeta_m)$ extends to $C$: extend the embedding to an
algebraic closure and restrict it to the normal extension $C$.
Thus $D_m$ is nonempty. It is closed, being the fiber of a continuous
map into a finite discrete group. Compatibility of character
coordinates gives $D_{m+1}\subseteq D_m$. The actual Krull Galois
group is compact, so these nested nonempty closed sets have a common
point $\sigma$. All reductions of $\chi(\sigma)$ and $u$ agree,
and a p-adic integer is determined by its reductions modulo every
$p^m$. Hence $\chi(\sigma)=u$. This proves surjectivity and,
together with the closed embedding, the topological group isomorphism.

If $C/\Qp$ had finite degree $N$, every subfield $E_m$ would have
degree at most $N$. Their degrees are $\varphi(p^m)=(p-1)p^{m-1}$
for $m\geq1$, which are unbounded since $p\geq2$. This contradicts
the bound $N$ and proves infinite degree, including at $p=2$.
\end{proof}

\begin{lemma}[The rational cyclotomic layers]
\label{ft:rational-cyclotomic}
\lean{FalseTateExtension.finrank_cyclotomicLayer_padic, FalseTateExtension.cyclotomicLayer_residueDegree, FalseTateExtension.irreducible_cyclotomicUniformizer, FalseTateExtension.cyclotomicLayer_ramificationIdx_eq_totient, FalseTateExtension.cyclotomicLayerGaloisEquivUnits, FalseTateExtension.cyclotomicExtensionGaloisEquivUnits, FalseTateExtension.not_finiteDimensional_cyclotomicExtension_padic}
\leanok
\source{serre-1979-local-fields:page-0085}
\dcref{Serre, Chapter IV, Section 4, Proposition 17}
\uses{ft:padic-cyclotomic-polynomial,ft:cyclotomic-total-ramification,ft:finite-cyclotomic-character,ft:infinite-cyclotomic-character}
For $n\geq1$, put $F_n=\Qp(\zeta_n)$ and
$d_n=(p-1)p^{n-1}$. Then $F_n/\Qp$ is totally ramified of degree $d_n$,
$\zeta_n-1$ is a uniformizer, and the cyclotomic character identifies
$\Gal(F_n/\Qp)$ with $(\mathbb Z/p^n\mathbb Z)^\times$.
Consequently, for $F_\infty=\bigcup_n F_n$,
$\Gal(F_\infty/\Qp)=\Zp^\times$ and $F_\infty/\Qp$ is infinite.
\end{lemma}
\begin{proof}
The polynomial
\[
 P_n(T)=\Phi_{p^n}(1+T)
       =\sum_{j=0}^{p-1}(1+T)^{jp^{n-1}}
\]
is monic of degree $d_n$ and has constant term $p$. In
$\mathbb F_p[T]$ it equals $T^{d_n}$: putting
$X=T^{p^{n-1}}$, its reduction is
$\sum_{j=0}^{p-1}(1+X)^j=X^{p-1}$.
The polynomial lemma gives $[F_n:\Qp]=d_n$.
The preceding total-ramification lemma gives residue degree one and
makes $\zeta_n-1$ a uniformizer.
All the primitive roots $\zeta_n^u$, for
$u\in(\mathbb Z/p^n\mathbb Z)^\times$, belong to $F_n$ and are the
distinct roots of the irreducible polynomial $\Phi_{p^n}$.
They give exactly the stated automorphisms. Taking inverse limits gives
$\Zp^\times$. The degrees $d_n$ are unbounded, including when $p=2$.
\end{proof}

\begin{lemma}[The finite cyclotomic residue-degree bound]
\label{ft:cyclotomic-finite-residue-bound}
\uses{ft:extensions,ft:padic-cyclotomic-polynomial,ft:cyclotomic-total-ramification,fw:local-dvr-degree-identity,fw:normalized-valuation-scaling,fw:ramification-field-equiv}
\lean{FalseTateExtension.padicUnitBall_residueFinite, FalseTateExtension.cyclotomicLayer_padic_upper_zero, FalseTateExtension.rationalCyclotomic_ramificationIdx, FalseTateExtension.finrank_cyclotomicLayer_le_totient, FalseTateExtension.rationalCyclotomicSubfield, FalseTateExtension.isPrimitiveRoot_cyclotomicLayer_generator, FalseTateExtension.rationalCyclotomicSubfield_isCyclotomicExtension, FalseTateExtension.rationalCyclotomicSubfield_finiteDimensional, FalseTateExtension.cyclotomicBase_unitBall_isDiscreteValuationRing, FalseTateExtension.cyclotomicLayer_spectral_residueDegree_le, FalseTateExtension.cyclotomicLayer_extensionValuationSubring_eq, FalseTateExtension.compatibleSpectralIntegerRingEquiv, FalseTateExtension.cyclotomicLayerIntegerRingEquiv, FalseTateExtension.coe_cyclotomicLayerIntegerRingEquiv, FalseTateExtension.cyclotomicLayerResidueEquiv, FalseTateExtension.cyclotomicLayerResidueEquiv_residue, FalseTateExtension.cyclotomicLayer_residueDegree_eq, FalseTateExtension.cyclotomicLayer_residueDegree_le}
\leanok
Let $K/\Qp$ be finite, with its compatible complete valuation,
and let $C_m=K(\zeta_m)$ for $m\geq1$, with its actual compatible
extension valuation. Then
\[
 f(C_m/K)\leq e(K/\Qp).
\]
The assertion holds for the actual chosen-root field and is
independent of the chosen compatible extension valuation.
\end{lemma}
\begin{proof}
Inside this same field $C_m$, put $F_m=\Qp(\zeta_m)$ and
$d_m=\varphi(p^m)$. The root remains primitive as an element of
$F_m$ and $C_m$, since all inclusions are injective. The actual
rational cyclotomic polynomial and total-ramification theorems,
transported along the $\Qp$-isomorphism carrying the original
chosen primitive root to this one, give
$[F_m:\Qp]=e(F_m/\Qp)=d_m$. Compatible extension valuations
have the same valuation ring, so this is the ramification index
of the actual inclusion $F_m\subseteq C_m$ under consideration.

The minimal polynomial of $\zeta_m$ over $K$ divides the image
of its rational minimal polynomial. Since $C_m=K(\zeta_m)$,
the degree of this minimal polynomial is $[C_m:K]$, whence
$[C_m:K]\leq d_m$.
Write $e_m=e(C_m/K)$, $f_m=f(C_m/K)$ and
$e_K=e(K/\Qp)$. The actual local degree identity and the two
ramification-index tower identities give
\[
 e_mf_m=[C_m:K]\leq d_m
 \leq d_m e(C_m/F_m)=e(C_m/\Qp)=e_K e_m.
\]
Here $e(C_m/F_m)\geq1$, and $e_m>0$ for a finite injective
local extension of DVRs. Cancelling this positive integer proves
$f_m\leq e_K$. Every index and residue degree in the argument
comes from the actual compatible valuation rings, so uniqueness
of those rings also gives the asserted independence.
\end{proof}

\begin{lemma}[Stabilization of cyclotomic inertia preimages]
\label{ft:cyclotomic-inertia-preimages}
\uses{ft:extensions,ft:cyclotomic-finite-residue-bound,fw:galois-ring-residue-theory,fw:local-inertia-size,fw:ordinary-upper-restriction,fw:infinite-upper-groups,fw:infinite-upper-restriction,fw:finite-base-residue}
\lean{FalseTateExtension.cyclotomicRestrictionToLayer, FalseTateExtension.cyclotomicRestrictionToLayer_commutes, FalseTateExtension.cyclotomicRestrictionToLayer_surjective, FalseTateExtension.continuous_cyclotomicRestrictionToLayer, FalseTateExtension.cyclotomicRestrictionBetweenLayers, FalseTateExtension.cyclotomicRestrictionBetweenLayers_commutes, FalseTateExtension.cyclotomicRestrictionToLayer_comp, FalseTateExtension.cyclotomicFiniteLayer, FalseTateExtension.cyclotomicFiniteLayerEquiv, FalseTateExtension.restrictNormal_cyclotomicFiniteLayer, FalseTateExtension.exists_le_cyclotomicFiniteLayer, FalseTateExtension.cyclotomicBase_residueFinite, FalseTateExtension.cyclotomicLayer_inertia_index, FalseTateExtension.cyclotomicInertiaPreimage, FalseTateExtension.cyclotomicInertiaPreimage_finiteIndex, FalseTateExtension.cyclotomicInertiaPreimage_normal, FalseTateExtension.isOpen_cyclotomicInertiaPreimage, FalseTateExtension.cyclotomicInertiaPreimage_index, FalseTateExtension.cyclotomicInertiaPreimage_index_le, FalseTateExtension.cyclotomicInertiaPreimage_antitone, FalseTateExtension.cyclotomicInertiaPreimage_index_monotone, FalseTateExtension.exists_cyclotomicInertiaPreimage_stable, FalseTateExtension.upperRamificationGroupInfinite_cyclotomic_eq_iInf, FalseTateExtension.exists_cyclotomic_inertia_eq_preimage, FalseTateExtension.cyclotomic_inertia_finiteIndex, FalseTateExtension.cyclotomic_inertia_index_le, FalseTateExtension.isOpen_cyclotomic_inertia, FalseTateExtension.cyclotomic_inertia_eq_absolute_image}
\leanok
Let $K/\Qp$ be finite, $C=K^{\mathrm{cyc}}$, and
$C_m=K(\zeta_m)$ for $m\geq1$. Write $D=\Gal(C/K)$,
$r_m:D\to\Gal(C_m/K)$ for ordinary restriction and
\[
 B_m=r_m^{-1}(\Gal(C_m/K)^0).
\]
Each $B_m$ is an open normal subgroup of finite index,
\[
 [D:B_m]=f(C_m/K)\leq e(K/\Qp),\qquad B_{m+1}\subseteq B_m.
\]
There exists $m_0\geq1$ such that
\[
 B_m=B_{m_0}\quad(m\geq m_0),\qquad
 D^0=\bigcap_{m\geq1}B_m=B_{m_0}.
\]
In particular the actual infinite upper group $D^0$ is open.
It is also the image of
$\Gal(\overline K/K)^0$ under ordinary restriction to $C$.
\end{lemma}
\begin{proof}
Each $C_m/K$ is finite Galois. Ordinary restriction is continuous
and surjective, because an automorphism of a normal intermediate
field extends to the algebraic closure and then preserves $C$.
Its target is finite discrete; the preimage of its normal inertia
subgroup is therefore open and normal. The finite residue extension
is Galois: normality follows from the actual residue reduction
homomorphism, and separability follows from perfection of the
base residue field. That residue field is finite, since $K/\Qp$
is finite and the residue field of $\Qp$ is $\mathbb F_p$.
The actual quotient by inertia is its residue automorphism group,
so its order is $f(C_m/K)$. Surjectivity of $r_m$ now gives the
index equality, and the finite residue-degree bound gives the
inequality. Finite upper restriction sends the next layer's zero
upper group into the previous layer's zero upper group; compatibility
of ordinary restriction proves $B_{m+1}\subseteq B_m$.

The positive integers $[D:B_m]$ form a nondecreasing sequence
bounded by $e(K/\Qp)$. Its range is finite and nonempty, so it
has a largest element, attained at some $m_0$. For $m\geq m_0$
the indices coincide. The inclusion $B_m\subseteq B_{m_0}$ and
the finite index formula
\[
 [D:B_m]=[D:B_{m_0}][B_{m_0}:B_m]
\]
force the relative index to be one. Thus $B_m=B_{m_0}$, and
monotonicity identifies their intersection with $B_{m_0}$.

Every finite intermediate field $E/K$ of $C/K$ is contained in
some $C_m$. Indeed choose a finite $K$-basis of $E$; each basis
element belongs to a layer because $C$ is their union. The largest
of these finitely many layer indices contains the basis and hence
all of $E$. For a finite Galois $E$, finite upper restriction
therefore shows that membership in every $B_m$ implies the
zero-upper condition on $E$. Conversely the infinite upper condition
implies every layer condition by its definition. This proves
$D^0=\bigcap_m B_m$. Its openness follows from stabilization.
Finally the proved infinite upper-restriction theorem, applied to
the actual inclusion $C\subseteq\overline K$ and $u=0$, identifies
$D^0$ with the image of $\Gal(\overline K/K)^0$.
This argument establishes openness directly from finite preimages;
it does not assert an identification with the infinite residue quotient.
\end{proof}

\begin{lemma}[The actual infinite cyclotomic residue bound]
\label{ft:cyclotomic-infinite-residue-bound}
\uses{ft:extensions,ft:cyclotomic-finite-residue-bound,fw:spectral-integer-residue-maps,fw:finite-base-residue}

\lean{FalseTateExtension.cyclotomicResidueToExtension, FalseTateExtension.cyclotomicResidueBetweenLayers, FalseTateExtension.cyclotomicResidueToExtension_comp, FalseTateExtension.exists_cyclotomicResidueToExtension_preimage, FalseTateExtension.cyclotomicResidue_degree_monotone, FalseTateExtension.exists_cyclotomicResidueToExtension_bijective, FalseTateExtension.exists_cyclotomicResidueEquiv, FalseTateExtension.cyclotomicExtension_residue_finiteDimensional, FalseTateExtension.cyclotomicExtension_residue_degree_le}
\leanok
Let $K/\Qp$ be finite and $C=K^{\mathrm{cyc}}$.
Let $\mathcal O_C$ be its actual spectral valuation ring and
$k_C$ its actual residue field. Then $k_C/k_K$ is finite with
\[
 [k_C:k_K]\leq e(K/\Qp).
\]
For some positive chosen finite layer $C_N$, the actual residue
map $k_{C_N}\to k_C$ is an isomorphism over the original $k_K$.
No discrete valuation or completeness assertion about $C$ is required.
\end{lemma}
\begin{proof}
Every inclusion of a chosen finite layer into a later layer or into
$C$ preserves its spectral norm. Indeed the spectral norm of an
algebraic extension of a complete nonarchimedean field is unique,
and its restriction along an injective algebraic field map is that
of the source. Thus the inclusions give actual local maps of their
valuation rings. The induced residue maps are field homomorphisms,
hence injective, and commute with one another and with the original
base residue map. They are therefore $k_K$-linear.

Every $z\in k_C$ has an integral representative $x\in\mathcal O_C$.
The underlying element $x\in C$ lies in a chosen finite layer,
since $C$ is their union. Norm preservation makes that representative
integral already in the finite layer; its residue maps to $z$.
Hence every element of the actual $k_C$ comes from a finite-layer
residue field.

Write $f_m=[k_{C_m}:k_K]$ for positive layers. Injectivity of the
actual residue maps makes $f_m$ nondecreasing; the finite-layer
bound gives $1\leq f_m\leq e(K/\Qp)$. Choose $N$ attaining the
largest value of this nonempty finite set of integer degrees.
For $m\geq N$, the injective $k_K$-linear map
$k_{C_N}\to k_{C_m}$ is between finite vector spaces of equal
dimension, so it is surjective. For any $z\in k_C$, choose a layer
$m$ representing it, pass to $\max(m,N)$, and use that surjectivity.
Commutativity of the actual residue maps shows that $z$ comes from
$k_{C_N}$. The map $k_{C_N}\to k_C$ is therefore bijective and
is an actual $k_K$-algebra isomorphism. Transporting finite
dimensionality along this map gives
$[k_C:k_K]=f_N\leq e(K/\Qp)$.
\end{proof}

\begin{lemma}[Inertia and residue fields]
\label{ft:inertia-residue}\notready
\source{serre-1979-local-fields:page-0069}
\dcref{Serre, Chapter IV, Section 1, Proposition 1}
Let $N/K$ be an algebraic Galois extension, and write $I_K$ for the
inertia subgroup of the absolute Galois group of $K$. The image $I_N$
of $I_K$ in $\Gal(N/K)$ is the kernel of reduction, and
\[
 \Gal(N/K)/I_N\simeq\Gal(k_N/k_K).
\]
This quotient is abelian. Inertia images commute with restriction to
Galois intermediate extensions. Moreover $I_N$ is open if and only if
$[k_N:k_K]$ is finite.
\end{lemma}
\begin{proof}
For a finite Galois extension $E/K$, the maximal unramified subfield
has residue field $k_E$ and degree $[k_E:k_K]$. To see its existence,
lift the irreducible polynomial of a generator of the finite field
extension $k_E/k_K$ to $\mathcal O_K[T]$. Hensel lifting a simple
residue root in $E$ gives a root of the lift; its field has that degree,
residue degree that degree, and therefore ramification index one.
The unramified field is unique in the fixed algebraic closure: its
residue elements have unique multiplicative lifts as roots of unity
of order prime to $p$, and those lifts generate it.
Every residue automorphism lifts to this subfield and extends to an
automorphism of $E$, since $E/K$ is normal. Thus reduction onto
$\Gal(k_E/k_K)$ is surjective. Its kernel is the inertia group, and
the image of absolute inertia is precisely this kernel: the maximal
unramified extension of $K$ has the union of all finite unramified
extensions as its field, so its intersection with $E$ is the subfield
just constructed.

Apply the finite statement to the finite Galois subextensions of $N/K$.
An element of $\Gal(N/K)$ acting trivially on $k_N$ restricts to inertia
at every finite level. Its lifts to $I_K$ form nonempty closed subsets
of the compact group $I_K$ with the finite intersection property
(use the compositum for finitely many levels). It therefore lifts to
$I_K$. The same compactness argument, lifting a prescribed residue
automorphism at finitely many levels, proves surjectivity of reduction.
The quotient is consequently the inverse limit of the cyclic Galois
groups of finite residue fields and is abelian. Restriction commutes
with images of the single group $I_K$. Finally a closed subgroup of a
profinite group is open exactly when it has finite index, which here
is equivalent to the residue extension being finite.
\end{proof}

\begin{lemma}[Naturality of the canonical cyclotomic character]
\label{ft:cyclotomic-character-naturality}
\lean{FalseTateExtension.cyclotomicCharacter_eq_of_intertwines}
\leanok
Let $E$ and $E'$ be characteristic-zero fields containing primitive
$p^n$-th roots of unity for every $n$, and let $i:E\hookrightarrow E'$
be a field embedding. If field automorphisms $g$ and $g'$ satisfy
$i\circ g=g'\circ i$, their canonical cyclotomic characters are equal.
\end{lemma}
\begin{proof}
Choose a primitive $p^n$-th root $\zeta\in E$. Its image $i(\zeta)$
is primitive of the same order, since the embedding is injective.
Write $b_n$ and $b'_n$ for the reductions of the two characters
modulo $p^n$. The defining character identities and the intertwining
identity give
\[
 i(\zeta)^{b_n}=i(g(\zeta))=g'(i(\zeta))=i(\zeta)^{b'_n}.
\]
Primitivity implies $b_n=b'_n$ in $\mathbb Z/p^n\mathbb Z$.
Thus the underlying p-adic integers have equal reductions for every
$n$, hence are equal. Units are determined by their underlying
elements, which proves equality of the characters. The argument
also covers $n=0$, where the residue ring has one element.
\end{proof}

\begin{lemma}[Openness of the cyclotomic image over a finite base field]
\label{ft:cyclotomic-image-open}
\uses{ft:extensions,ft:infinite-cyclotomic-character,ft:cyclotomic-projection-image,ft:cyclotomic-character-naturality}
\lean{FalseTateExtension.cyclotomicCharacter_algebraicClosure_padic_surjective, FalseTateExtension.cyclotomicCharacter_surjective_of_isAlgClosure, FalseTateExtension.range_restrictScalarsHom_eq_fixingSubgroup, FalseTateExtension.index_range_restrictScalarsHom, FalseTateExtension.cyclotomicCharacter_restrict_cyclotomicExtension, FalseTateExtension.cyclotomicExtensionCharacter_range_eq_map, FalseTateExtension.cyclotomicExtensionCharacter_finiteIndex, FalseTateExtension.cyclotomicExtensionCharacter_index_dvd, FalseTateExtension.cyclotomicExtensionCharacter_index_pos, FalseTateExtension.cyclotomicExtensionCharacter_index_bounds, FalseTateExtension.isOpen_cyclotomicExtensionCharacter_range}
\leanok
Let $K/\Qp$ be finite and let $C=K(\mu_{p^\infty})$ be the
actual cyclotomic extension defined by the chosen roots. Its
canonical character image $H\subseteq\Zp^\times$ has finite index,
\[
 [\Zp^\times:H]\mid[K:\Qp],\qquad
 1\leq[\Zp^\times:H]\leq[K:\Qp],
\]
and is open. This holds at every prime.
\end{lemma}
\begin{proof}
Put $\Omega=\overline K$, with its composite $\Qp$-algebra
structure. Since $K/\Qp$ is algebraic, $\Omega$ is also an
algebraic closure of $\Qp$. Choose a $\Qp$-isomorphism from
$\overline{\Qp}$ to $\Omega$. On $\overline{\Qp}$, restriction
to the rational cyclotomic extension is surjective: every embedding
of that normal extension extends to the algebraic closure. Its
character is surjective by the infinite rational character lemma.
Naturality therefore proves that the canonical character
\[
 \chi_\Omega:\Gal(\Omega/\Qp)\longrightarrow\Zp^\times
\]
is surjective as well.

Let $\kappa$ be the actual image of $K$ in $\Omega$, and let
$S\subseteq\Gal(\Omega/\Qp)$ be the scalar-restriction image
of $\Gal(\Omega/K)$. Then $S$ is exactly the subgroup fixing
$\kappa$: the forward inclusion follows from $K$-linearity;
an automorphism fixing $\kappa$ is $K$-linear and gives the reverse
inclusion. Galois correspondence gives
$[\Gal(\Omega/\Qp):S]=[\kappa:\Qp]=[K:\Qp]$.
In particular $S$ has finite index.
Ordinary restriction $\Gal(\Omega/K)\to\Gal(C/K)$ is
surjective, and naturality of the character shows
$H=\chi_\Omega(S)$. For a surjective group homomorphism,
the cosets of a subgroup map surjectively to the cosets of its
image; the image index divides the original finite index, by
the subgroup index formula applied to the preimage of the image.
Consequently $H$ has finite index dividing $[K:\Qp]$.
Both indices are positive integers, so the displayed bounds follow.
Finally the actual character of $C/K$ is a closed embedding;
its image $H$ is closed. A closed subgroup of finite index in
a profinite group is open, proving the assertion.
\end{proof}

\begin{lemma}[Cyclotomic image and cyclotomic inertia]
\label{ft:cyclotomic-inertia}
\uses{ft:extensions,ft:rational-cyclotomic,ft:cyclotomic-image-open,ft:cyclotomic-finite-residue-bound,ft:cyclotomic-inertia-preimages,ft:cyclotomic-infinite-residue-bound}
\lean{FalseTateExtension.cyclotomicInertiaCharacterImage, FalseTateExtension.cyclotomicInertiaCharacterImage_eq_map, FalseTateExtension.cyclotomicInertiaCharacterImage_eq_absolute_character, FalseTateExtension.isOpenEmbedding_cyclotomicExtensionCharacter, FalseTateExtension.isOpen_cyclotomicInertiaCharacterImage, FalseTateExtension.cyclotomicInertiaCharacterImage_finiteIndex, FalseTateExtension.padicUnits_infinite, FalseTateExtension.exists_ne_one_mem_cyclotomicInertiaCharacterImage}
\leanok
Put $C=K^{\mathrm{cyc}}$ and
$H=\Gal(C/K)\subset\Zp^\times$. Both $H$ and the image
$J=\operatorname{im}(I_K\longrightarrow H)$ are open in
$\Zp^\times$. In particular $J$ contains an element $h\ne1$.
Furthermore $[k_C:k_K]\leq e(K/\Qp)$.
\end{lemma}
\begin{proof}
The actual character of $C/K$ is a closed embedding with open image
$H$, by the cyclotomic image theorem. It is therefore an open
embedding: it is a homeomorphism onto its image, and that image is
open in $\Zp^\times$.
The actual infinite zero upper group $\Gal(C/K)^0$ is open by
finite inertia-preimage stabilization. Its image under this open
embedding is open in $\Zp^\times$. Ordinary infinite upper
restriction identifies it with the stated image of absolute inertia,
so this image is precisely $J$.

The group $\Zp^\times$ is infinite: the distinct elements
$1+pn$ for $n\geq0$ are units, since their reduction modulo $p$
is one, and they are distinct by characteristic zero.
An open subgroup of this compact group has finite index: its
coset space is discrete and compact, hence finite. If that subgroup
were finite, the group would be a finite union of finite cosets,
a contradiction. Thus $J$ is infinite and has an element different
from one. The actual infinite residue bound just proved supplies
$[k_C:k_K]\leq e(K/\Qp)$ independently of any identification
of an infinite residue quotient.
\end{proof}

\begin{lemma}[Compatibility of finite coordinates under restriction]
\label{ft:coordinate-transitions}
\source{liu-2008-lattices:page-0022}
\dcref{Liu, Section 5.1; transition maps on compatible roots}
\uses{ft:finite-coordinates,ft:extensions}
\lean{FalseTateExtension.finiteGaloisLayer, FalseTateExtension.finiteLayerRestriction, FalseTateExtension.finiteLayerRestriction_commutes, FalseTateExtension.finiteLayerRestriction_surjective, FalseTateExtension.finiteLayerRestriction_refl, FalseTateExtension.finiteLayerRestriction_comp, FalseTateExtension.radical_add_pow, FalseTateExtension.radical_pow_sub, FalseTateExtension.rootOfUnity_add_pow, FalseTateExtension.rootOfUnity_pow_sub, FalseTateExtension.inclusion_layerRadical, FalseTateExtension.inclusion_layerRootOfUnity, FalseTateExtension.finiteLayerRestriction_radical, FalseTateExtension.finiteLayerRestriction_rootOfUnity, FalseTateExtension.layerRootOfUnity_pow_reduction, FalseTateExtension.finiteKummerExponent_restrict, FalseTateExtension.finiteCyclotomicCharacter_restrict}
\leanok
Let $\pi\ne0$ and $n\leq m$. Put $E_j=K(\pi_j,\zeta_j)$ and
let $r_{m,n}:\Gal(E_m/K)\to\Gal(E_n/K)$ be the ordinary finite
Galois restriction. Let $\rho_{m,n}:\mathbb Z/p^m\mathbb Z\to
\mathbb Z/p^n\mathbb Z$ be reduction and $\rho_{m,n}^{\times}$ its
map on units. Then
\[
 b_n(r_{m,n}(g))=\rho_{m,n}(b_m(g)),\qquad
 \chi_n(r_{m,n}(g))=\rho_{m,n}^{\times}(\chi_m(g)).
\]
The restrictions are the identity for $m=n$ and compose in towers.
\end{lemma}
\begin{proof}
Iterating the chosen root compatibilities gives
$\pi_n=\pi_m^{p^{m-n}}$ and $\zeta_n=\zeta_m^{p^{m-n}}$.
Apply $g$ to the first identity and substitute its radical evaluation
at level $m$. Commuting multiplication and powers gives
\[
 r_{m,n}(g)(\pi_n)=\zeta_n^{b_m(g)}\pi_n.
\]
Powers of $\zeta_n$ depend only on the exponent modulo $p^n$.
Uniqueness of the radical coordinate therefore gives the first formula.
The same calculation with the root-of-unity identity gives
$r_{m,n}(g)(\zeta_n)=\zeta_n^{\chi_m(g)}$. Reduction preserves units,
and character uniqueness gives the second formula. Ordinary restriction
preserves values on the smaller field, so its identity and composition
laws follow by evaluation there. At $n=0$ the target coordinate groups
are trivial, and the same formulas apply.
\end{proof}

\begin{lemma}[Actual $p$-adic cyclotomic and Kummer coordinates]
\label{ft:padic-coordinates}
\source{liu-2008-lattices:page-0022}
\dcref{Liu, Section 5.1; compatible-coordinate construction over a characteristic-zero base}
\uses{ft:coordinate-transitions,ft:galois-layers,ft:extensions}
\lean{FalseTateExtension.infiniteLayerRestriction, FalseTateExtension.infiniteLayerRestriction_commutes, FalseTateExtension.infiniteLayerRestriction_surjective, FalseTateExtension.infiniteLayerRestriction_comp, FalseTateExtension.infiniteLayerRestriction_jointly_injective, FalseTateExtension.extensionRootOfUnity, FalseTateExtension.coe_extensionRootOfUnity, FalseTateExtension.isPrimitiveRoot_extensionRootOfUnity, FalseTateExtension.falseTateExtension_hasEnoughRootsOfUnity, FalseTateExtension.padicCyclotomicCharacter, FalseTateExtension.padicCyclotomicCharacter_toZModPow, FalseTateExtension.kummerRepresentative, FalseTateExtension.kummerRepresentative_step, FalseTateExtension.padicKummerExponent, FalseTateExtension.padicKummerExponent_toZModPow, FalseTateExtension.padicKummerExponent_unique, FalseTateExtension.padicKummerExponent_mul, FalseTateExtension.padicCoordinates_injective}
\leanok
Let $\pi\ne0$ and $L=K^{\mathrm{FT}}_\infty$. For every $n$, ordinary
restriction $r_n:\Gal(L/K)\to\Gal(E_n/K)$ is surjective, commutes
with the original field inclusion, and satisfies
$r_{m,n}r_m=r_n$. The actual field $L$ contains primitive $p^n$-th
roots of unity for every $n$. Its usual cyclotomic character restricts
to a homomorphism $\chi:\Gal(L/K)\to\Zp^\times$ with
\[
 \chi(g)\bmod p^n=\chi_n(r_n(g)).
\]
There is a unique map $b:\Gal(L/K)\to\Zp$ satisfying
\[
 b(g)\bmod p^n=b_n(r_n(g)).
\]
It satisfies $b(gh)=b(g)+\chi(g)b(h)$, and
$g\mapsto(b(g),\chi(g))$ is injective, for every prime $p$.
\end{lemma}
\begin{proof}
Each $E_n/K$ and $L/K$ is actually Galois. The normal-extension
embedding theorem extends every automorphism of $E_n$ to $L$, proving
surjectivity of ordinary restriction. Evaluation in the ambient field
proves its inclusion and transition identities.

The element $\zeta_n$ belongs to $E_n\subset L$ and retains its
primitive order under this injective inclusion. Thus $L$ has enough
$p^n$-th roots of unity for every $n$. The cyclotomic character on
$L$ is consequently characterized by its action on these actual roots.
Its projection at $p^n$ agrees with $\chi_n(r_n(g))$ by the finite
character evaluation and uniqueness theorem.

The finite transition lemma makes $b_n(r_n(g))$ compatible under all
reductions. Choose integer representatives $a_n$. For $m\geq n$,
the difference $a_m-a_n$ is divisible by $p^n$. Hence $(a_n)$ is a
$p$-adic Cauchy sequence and converges to an element $b(g)\in\Zp$.
Its projection modulo $p^n$ is the prescribed finite value. Changing
representatives changes none of these projections, and an element of
$\Zp$ is uniquely determined by all its projections, so this defines
the unique stated map.

Project $b(gh)$ and $b(g)+\chi(g)b(h)$ modulo every $p^n$.
Restriction is a homomorphism, the cyclotomic projection has just been
identified, and the finite cocycle identity identifies these
projections. Uniqueness of $p$-adic projections proves the cocycle law.
Equal pairs give equal finite pairs at every level, so finite-coordinate
injectivity identifies the restrictions on every $E_n$. Every element
of $L$ lies in a diagonal layer, and the original field inclusions are
injective. The two automorphisms therefore agree on all of $L$.
\end{proof}

\begin{lemma}[The actual closed topological affine embedding]
\label{ft:affine-closed-embedding}
\source{liu-2008-lattices:page-0022}
\dcref{Liu, Section 5.1; topology of the compatible-coordinate realization}
\uses{ft:padic-coordinates,ft:galois-layers}
\lean{FalseTateExtension.padicAffineAction, FalseTateExtension.PadicAffineGroup, FalseTateExtension.padicAffineEquivProd, FalseTateExtension.padicAffineTopologicalSpace, FalseTateExtension.padicAffineHomeomorphProd, FalseTateExtension.padicAffineT2Space, FalseTateExtension.continuous_padicAffine_left, FalseTateExtension.continuous_padicAffine_right, FalseTateExtension.padicAffineIsTopologicalGroup, PadicInt.continuous_of_toZModPow, FalseTateExtension.continuous_infiniteLayerRestriction, FalseTateExtension.continuous_padicKummerExponent, FalseTateExtension.continuous_padicCyclotomicCharacter, FalseTateExtension.affineGaloisHom, FalseTateExtension.affineGaloisHom_coordinates, FalseTateExtension.affineGaloisHom_injective, FalseTateExtension.continuous_affineGaloisHom, FalseTateExtension.isClosedEmbedding_affineGaloisHom}
\leanok
For every prime $p$ and $\pi\ne0$, give the usual semidirect product
$\Zp\rtimes\Zp^\times$, with action $u\cdot b=ub$, its product
topology. It is a Hausdorff topological group. The actual map
\[
 \Gal(K^{\mathrm{FT}}_\infty/K)\longrightarrow\Zp\rtimes\Zp^\times,
 \qquad g\longmapsto(b(g),\chi(g))
\]
is a continuous group homomorphism and a closed topological embedding.
The source has its actual Krull topology.
\end{lemma}
\begin{proof}
The group law is $(b,u)(c,v)=(b+uc,uv)$, and inversion is
$(b,u)^{-1}=(-u^{-1}b,u^{-1})$. Continuity of addition and
multiplication on $\Zp$ and of the unit operations proves the
topological group assertion for the product topology; both factors are
Hausdorff. The preceding coordinate theorem proves the homomorphism
law and injectivity.

Ordinary restriction to an actual finite Galois layer is continuous:
the preimage of its identity is the open subgroup fixing that layer,
and the fibers are its cosets. In particular every finite projection
of $b$ factors through a continuous map to a finite discrete Galois
group and is continuous. Residue cosets modulo $p^n$ form a basis of
the actual $p$-adic topology on $\Zp$. Indeed their kernels are
$p^n\Zp$, and their translates are the closed balls of radius $p^{-n}$;
these balls are open as well. Given a positive real radius, choose $n$
with $p^{-n}$ smaller than that radius. Thus every neighborhood contains
such a residue coset. Continuity of all finite projections therefore
proves continuity of $b$. The cyclotomic coordinate is the usual
continuous cyclotomic character restricted to the actual $K$-algebra
automorphisms. Hence the displayed map is continuous in the product
topology.

The actual infinite false Tate extension is Galois, so its Krull Galois
group is compact. A continuous injection from a compact space to a
Hausdorff space is a homeomorphism onto its image. The image is compact
and therefore closed. These facts give the claimed closed embedding.
\end{proof}

\begin{lemma}[The actual cyclotomic projection image]
\label{ft:cyclotomic-projection-image}
\source{liu-2008-lattices:page-0022}
\dcref{Liu, Section 5.1; the cyclotomic projection}
\uses{ft:affine-closed-embedding,ft:extensions,ft:galois-layers}
\lean{FalseTateExtension.cyclotomicRootOfUnity, FalseTateExtension.isPrimitiveRoot_cyclotomicRootOfUnity, FalseTateExtension.cyclotomicExtension_hasEnoughRootsOfUnity, FalseTateExtension.cyclotomicExtensionCharacter, FalseTateExtension.continuous_cyclotomicExtensionCharacter, FalseTateExtension.cyclotomicExtensionCharacter_spec, FalseTateExtension.cyclotomicExtensionCharacter_injective, FalseTateExtension.isClosedEmbedding_cyclotomicExtensionCharacter, FalseTateExtension.cyclotomicExtension_le_falseTateExtension, FalseTateExtension.cyclotomicRestriction, FalseTateExtension.cyclotomicRestriction_commutes, FalseTateExtension.cyclotomicRestriction_surjective, FalseTateExtension.padicCyclotomicCharacter_eq_comp, FalseTateExtension.continuous_cyclotomicRestriction, FalseTateExtension.padicCyclotomicCharacter_ker, FalseTateExtension.padicCyclotomicCharacter_range, FalseTateExtension.cyclotomicImageEquiv, FalseTateExtension.cyclotomicImageEquiv_apply, FalseTateExtension.cyclotomicImageEquiv_restrict}
\leanok
Put $C=K^{\mathrm{cyc}}$, $L=K^{\mathrm{FT}}_\infty$, and
$G=\Gal(L/K)$. Restriction $r:G\to\Gal(C/K)$ is a continuous
surjection. The canonical cyclotomic character $\chi_C$ is injective
and continuous, and $\chi=\chi_C\circ r$. Therefore
$\ker\chi=\ker r$, and the actual projection image $\chi(G)$ is
isomorphic as a topological group to $\Gal(C/K)$.
\end{lemma}
\begin{proof}
Each finite cyclotomic layer is the splitting field of $X^{p^n}-1$;
their increasing union $C$ is Galois over $K$, as is $L$. The usual
extension theorem for embeddings into a normal algebraic extension
makes restriction onto $C$ surjective. Equivalently, extend an
embedding of $C$ to the algebraic closure and restrict the extension to
$L$, which is stable by normality. The continuity of restriction will follow from the actual character
factorization below.

The chosen primitive root $\zeta_n$ belongs to $C$ for every $n$,
so this field contains all $p$-power roots of unity. The canonical
cyclotomic character describes their actual action. If two
automorphisms have the same character, they agree on every $\zeta_n$,
then on each field $K(\zeta_n)$ and hence on its union $C$. This proves
injectivity. For $g\in G$, restriction and the original automorphism
act identically on every $\zeta_n$. Their character residues modulo
$p^n$ therefore agree for all $n$, and separation of $\Zp$ gives
$\chi(g)=\chi_C(r(g))$. Injectivity gives equality of the kernels,
and surjectivity gives $\chi(G)=\chi_C(\Gal(C/K))$.
Finally $\Gal(C/K)$ is compact in its Krull topology, the unit group
$\Zp^\times$ is Hausdorff, and the canonical character is continuous.
The injective map $\chi_C$ is therefore a closed embedding, so it
induces the actual topology on its domain. The identity
$\chi=\chi_C\circ r$ and continuity of $\chi$ now imply continuity
of $r$ through this inducing topology. The same closed embedding is a
homeomorphism onto its image. Since it is a group homomorphism, this
gives the stated topological group isomorphism.
\end{proof}

\begin{lemma}[The actual closed additive affine kernel]
\label{ft:affine-additive-kernel}
\source{liu-2008-lattices:page-0022}
\dcref{Liu, Section 5.1; the Kummer kernel in affine coordinates}
\uses{ft:affine-closed-embedding}
\lean{FalseTateExtension.continuous_padicAffine_inl, FalseTateExtension.affineAdditiveKernel, FalseTateExtension.mem_affineAdditiveKernel_iff, FalseTateExtension.padicKummerExponent_mem_affineAdditiveKernel, FalseTateExtension.kummerKernelHom, FalseTateExtension.kummerKernelHom_apply, FalseTateExtension.kummerKernelHom_bijective, FalseTateExtension.continuous_kummerKernelHom, FalseTateExtension.kummerKernelEquiv, FalseTateExtension.kummerKernelEquiv_apply}
\leanok
Let $f:G\to\Zp\rtimes\Zp^\times$ be the actual closed affine
embedding, $H=f(G)$, and let $\chi$ be its right coordinate. Then
\[
 M=\{a\in\Zp:(a,1)\in H\}
\]
is a closed additive subgroup. The actual Kummer coordinate gives a
topological group isomorphism $\ker\chi\longrightarrow M$, viewing the
target additively. Equivalently this is an additive topological group
isomorphism from the additive type tag of $\ker\chi$ onto $M$.
\end{lemma}
\begin{proof}
The canonical inclusion $\iota:a\mapsto(a,1)$ from the additive
$p$-adic group into the semidirect product is an injective continuous
group homomorphism. Its image is exactly the kernel of the right
projection. The inverse image of the closed subgroup $H$ under $\iota$
therefore gives the displayed closed additive subgroup $M$.

For $g\in\ker\chi$, the actual affine coordinate identity says
$f(g)=(b(g),1)$, so $b(g)\in M$. The cocycle identity reduces on this
kernel to $b(gh)=b(g)+b(h)$. Injectivity of $f$ makes this map
injective. If $a\in M$, by definition $(a,1)=f(g)$ for some $g\in G$;
comparison of right coordinates gives $g\in\ker\chi$ and comparison of
left coordinates gives $b(g)=a$. This proves surjectivity and the
actual coordinate formula.

The kernel is closed in the compact Krull group, since it is the
preimage of $\{1\}$ under the continuous cyclotomic character. Thus it
is compact. The coordinate map is continuous by restriction of the
actual continuous Kummer coordinate, while $M$ is Hausdorff as a
subgroup of $\Zp$. The continuous bijection is a homeomorphism, which
proves the asserted topological identification.
\end{proof}

\begin{lemma}[A closed affine realization]
\label{ft:affine-embedding}\uses{ft:extensions,ft:affine-closed-embedding,ft:cyclotomic-projection-image,ft:affine-additive-kernel}
\lean{FalseTateExtension.falseTateExtension_isGalois, FalseTateExtension.affineGaloisHom, FalseTateExtension.affineGaloisHom_coordinates, FalseTateExtension.isClosedEmbedding_affineGaloisHom, FalseTateExtension.cyclotomicImageEquiv, FalseTateExtension.cyclotomicImageEquiv_restrict, FalseTateExtension.padicCyclotomicCharacter_ker, FalseTateExtension.padicCyclotomicCharacter_range, FalseTateExtension.affineAdditiveKernel, FalseTateExtension.kummerKernelEquiv, FalseTateExtension.kummerKernelEquiv_apply}
\leanok
\source{liu-2008-lattices:page-0022}
\dcref{Liu, Section 5.1, action on compatible radicals}
Put $L=K^{\mathrm{FT}}_\infty$ and $G=\Gal(L/K)$. For every prime $p$,
$L/K$ is Galois, and there is a closed topological embedding
\[
 G\lhook\joinrel\longrightarrow\Zp\rtimes\Zp^\times,
 \qquad g\longmapsto(b(g),\chi(g)),
\]
where the multiplication is $(b,u)(c,v)=(b+uc,uv)$.
The projection image is $H=\Gal(K^{\mathrm{cyc}}/K)$, and its kernel
identifies with a closed additive subgroup $M$ of $\Zp$.
\end{lemma}
\begin{proof}
The preceding closed-embedding lemma supplies the displayed actual
homomorphism and its topological conclusion. Restriction onto the
cyclotomic Galois subextension is surjective. Its kernel is precisely
the intersection of this image with $\Zp\times\{1\}$ and is therefore
a closed additive subgroup. This construction requires no splitting
of the image as a semidirect product.
\end{proof}

\begin{lemma}[Closed additive subgroups of the p-adic integers]
\label{ft:padic-closed-additive-subgroups}
\source{liu-2008-lattices:page-0022}
\dcref{Liu, Section 5.1; closed translation subgroups}
\lean{PadicInt.mul_mem_closedAddSubgroup, PadicInt.closedAddSubgroupIdeal, PadicInt.mem_closedAddSubgroupIdeal, PadicInt.closedAddSubgroupIdeal_toAddSubgroup, PadicInt.closedAddSubgroup_eq_bot_or_span_pow, PadicInt.closedAddSubgroup_map_mulLeft_le, PadicInt.closedAddSubgroupIdeal_map_mulLeft_toAddSubgroup, PadicInt.isOpen_closedAddSubgroup_map_mulLeft, PadicInt.closedAddSubgroup_map_mulLeft_finiteRelIndex, PadicInt.closedAddSubgroup_map_mulLeft_relIndex_pos, PadicInt.closedAddSubgroup_map_mulLeft_relIndex_eq_one_of_eq_bot}
\leanok
For every prime $p$, each closed additive subgroup $M\subset\Zp$ is
zero or $p^s\Zp$ for some $s\geq0$. For every nonzero $a\in\Zp$,
the actual image $aM$ has positive finite index in $M$; if $M=0$,
this index is one.
\end{lemma}
\begin{proof}
For $x\in M$, all natural integer multiples belong to $M$. The natural
integers are dense in $\Zp$, multiplication by $x$ is continuous, and
$M$ is closed. Thus $ax\in M$ for every $a\in\Zp$. Consequently the
same carrier $M$ is an ideal of the discrete valuation ring $\Zp$.
Its ideals are zero or the powers $(p^s)$ of the maximal ideal,
which proves the classification.

If $M\ne0$ and $a\ne0$, the image $aM$ is the nonzero ideal
$(ap^s)$, contained in $M$. Every nonzero ideal of this compact
discrete valuation ring is open. Thus $aM$ is open in the compact
subgroup $M$, and its coset quotient is discrete and compact, hence
finite. Its index is positive because the quotient is nonempty.
When $M=0$, multiplication has the same zero image and the quotient
has one element, giving index one.
\end{proof}

\begin{lemma}[Closed additive subgroups and commutators]
\label{ft:affine-commutator}\uses{ft:padic-closed-additive-subgroups}
\lean{FalseTateExtension.padicAffine_commutator_inl, FalseTateExtension.padicAffineCompactSpace, FalseTateExtension.affineProjection, FalseTateExtension.continuous_affineProjection, FalseTateExtension.affineTranslationKernel, FalseTateExtension.affineTranslationHom, FalseTateExtension.affineTranslationHom_injective, FalseTateExtension.affineTranslationHom_range, FalseTateExtension.affineProjection_index, FalseTateExtension.affineTranslation_index, FalseTateExtension.affineTranslation_index_product, FalseTateExtension.affineTranslation_commutator, FalseTateExtension.affineTranslation_scalar_mem, FalseTateExtension.affineTranslation_scalar_image_le, FalseTateExtension.affineProjection_kernel_finiteRelIndex, FalseTateExtension.affineProjection_finiteIndex, FalseTateExtension.isOpen_of_affineProjection}
\leanok
Let $G\subset\Zp\rtimes\Zp^\times$ be a closed subgroup with projection
image $H$ and additive kernel $M$. Let $I\triangleleft G$ be closed, and assume its projection image
$J\subset H$ is open and contains $h\ne1$. Then $I$ is open in $G$.
More precisely,
\[
 (h-1)M\subset [G,G]\cap I\cap M,
 \qquad [G:I]=[H:J][M:M\cap I]<\infty.
\]
\end{lemma}
\begin{proof}
The preceding closed-additive-subgroup lemma makes multiplication by
the nonzero scalar $h-1$ have finite positive index in $M$, including
index one if $M=0$. Choose $(b,h)\in I$ above $h$.
The inverse in the affine group is $(b,h)^{-1}=(-h^{-1}b,h^{-1})$,
and direct multiplication gives, for every $m\in M$,
\[
 [(b,h),(m,1)]=((h-1)m,1).
\]
Every such commutator lies in $[G,G]$, in $M$, and in $I$ (also by
normality of $I$). Thus $M\cap I$ has finite index in $M$.
The preimage of $J$ in $G$ is $IM$: given an element over $J$, choose
an element of $I$ with the same projection and divide by it.
The quotient isomorphisms
$G/IM\simeq H/J$ and $IM/I\simeq M/(M\cap I)$ now prove the index
formula. The projection image $H$ is compact: the affine group has compact
$p$-adic and unit factors, and $G$ is closed. Thus the open subgroup
$J$ has finite index in $H$. The subgroup $I$ is closed of finite index,
so its finitely many closed cosets have open complements; in particular
$I$ is open.
\end{proof}

\begin{lemma}[Finite residue extension and Lie group structure]
\label{ft:lie-structure}
\uses{ft:affine-embedding,ft:cyclotomic-inertia,ft:affine-commutator,ft:inertia-residue}
\notready
For every prime $p$, $K^{\mathrm{FT}}_\infty/K$ is an infinite Galois
extension with finite residue extension and a $p$-adic Lie Galois group.
The same assertions hold for $K^{\mathrm{cyc}}/K$.
\end{lemma}
\begin{proof}
Use the closed affine realization and let $I$ be the image of $I_K$ in
$G$. By the residue lemma $I$ is closed and normal and $G/I$ is abelian.
Its image in $H$ is exactly the cyclotomic inertia image $J$, which is
open and contains $h\ne1$. The commutator lemma shows that $I$ is open;
the residue lemma then gives $[k_L:k_K]<\infty$.
The cyclotomic assertion on residue fields was already proved above.

For completeness the Lie group assertion also follows from this affine
realization. The group $H$ is open in $\Zp^\times$. If $M=0$, projection
is a continuous bijection $G\longrightarrow H$ between compact Hausdorff
groups and is a topological isomorphism; transport the $p$-adic Lie
structure of $H$ to $G$. If $M\ne0$, the preceding proof gives
$M=p^s\Zp$. Every left coset of $G$ in $\Zp\rtimes H$ has a
representative $(b,1)$, by multiplying by an element of $G$ with the same
projection. Two such representatives give the same coset exactly when
their difference lies in $M$. Thus the index is $[\Zp:M]=p^s$.
The subgroup $G$ is closed and of finite index, hence open in
$\Zp\rtimes H$. The latter is a $p$-adic Lie group: its two coordinates
have analytic multiplication $(b,u)(c,v)=(b+uc,uv)$ and inverse
$(-u^{-1}b,u^{-1})$. Its open subgroup $G$ is therefore a $p$-adic Lie
group. Finally both extensions contain the cyclotomic tower, whose
degrees over $K$ are unbounded because
$[KF_n:K]\geq d_n/[K:\Qp]$. They are infinite.
\end{proof}

\begin{remark}[Relation to the odd-prime splitting]
The affine coordinates are those underlying Liu\textquotesingle s description
of the false Tate Galois group. The argument above only uses closedness
of their image and thus applies at $p=2$ as well. It proves finite
residue degree over general $K$; total ramification is a separate
assertion over $\Qp$ in the odd-prime setting below.
\end{remark}

\subsection{A direct all-prime strict APF argument}

The following argument uses the actual radical tower and affine action.
It allows a finite initial residue extension in the false Tate tower.

\begin{lemma}[The actual finite radical layers and their differents]
\label{ft:radical-layer-ramification}
\uses{ft:extensions,fw:totally-ramified-uniformizer-generation,fw:monogenic-different,fw:norm-valuation,fw:local-dvr-degree-identity,fw:normalized-valuation-scaling,fw:different-order-transitivity}
\lean{FalseTateExtension.X_pow_sub_uniformizer_isEisensteinAt, FalseTateExtension.irreducible_X_pow_sub_uniformizer, FalseTateExtension.minpoly_root_uniformizer, FalseTateExtension.radicalLayerRoot, FalseTateExtension.radicalLayerRoot_pow, FalseTateExtension.finrank_radicalLayer, FalseTateExtension.minpoly_radicalLayerRoot, FalseTateExtension.norm_radicalLayerRoot, FalseTateExtension.radicalLayer_normalizedValuation_pow, FalseTateExtension.radicalLayer_residueDegree, FalseTateExtension.radicalLayer_normalizedValuation_root, FalseTateExtension.radicalLayerRoot_mem, FalseTateExtension.radicalUniformizer, FalseTateExtension.irreducible_radicalUniformizer, FalseTateExtension.radicalLayer_ramificationIdx, FalseTateExtension.adjoin_radicalUniformizer_eq_top, FalseTateExtension.minpoly_radicalUniformizer, FalseTateExtension.radicalLayer_different_order, FalseTateExtension.radicalLayer_different_order_toNat, FalseTateExtension.radicalLayer_step_finrank, FalseTateExtension.radicalLayer_step_ramificationIdx, FalseTateExtension.radicalLayer_step_residueDegree, FalseTateExtension.radicalLayer_step_different_order, FalseTateExtension.addVal_p_eq_padic_ramificationIdx, FalseTateExtension.addVal_p_toNat_eq_padic_ramificationIdx}
\leanok
Let $K$ be a complete discretely valued field of characteristic zero,
let $p$ be prime and let $\pi$ be an actual uniformizer of $K$.
For compatible $p$-power roots put $F_n=K(\pi_n)$ with
$\pi_n^{p^n}=\pi$. For every $n\geq0$, the original field $F_n$
has degree $p^n$ and minimal polynomial $X^{p^n}-\pi$, with
\[
 N_{F_n/K}(\pi_n)=(-1)^{p^n+1}\pi.
\]
For every compatible extension valuation its actual residue degree
is one, its ramification index is $p^n$, and $\pi_n$ is an actual
integral uniformizer generating the integer ring. Its actual different
exponent is
\[
 d_n=p^n n\,v_K(p)+p^n-1.
\]
Consecutive original radical fields have relative degree and
ramification index $p$, residue degree one, and actual different
exponent $p^{n+1}v_K(p)+p-1$. If the norm on $K$ is compatible
with its $\mathbb Q_p$-algebra structure, then $v_K(p)$ equals
its actual ramification index over $\mathbb Q_p$.
\end{lemma}
\begin{proof}
For every positive $d$, the polynomial $X^d-\pi$ is Eisenstein
at the actual maximal ideal: its leading coefficient is a unit,
all lower coefficients lie in that ideal, and its constant term
is outside its square. It is therefore irreducible over $K$.
Apply this with $d=p^n$ to the actual chosen root. Its power basis
computes both the degree and minimal polynomial, including $n=0$;
the constant coefficient gives the displayed signed norm.
For any compatible extension valuation, the normalized norm identity
gives $f v_{F_n}(\pi_n)=1$. Both are integers and $f>0$, so
$f=1$ and $v_{F_n}(\pi_n)=1$. The fundamental degree identity
then gives $e=p^n$. Thus the root is integral and is an actual
uniformizer. Total ramification makes its integral power basis
generate the actual integer ring. The monogenic different is
therefore generated by $p^n\pi_n^{p^n-1}$. Its normalized
valuation is $p^n v_K(p^n)+p^n-1=d_n$.

Absolute field-degree and ramification-index transitivity give
relative degree and ramification index $p$ at each consecutive
step. The relative degree identity then forces relative residue
degree one. Different transitivity gives
$d_{n+1}=d(F_{n+1}/F_n)+p d_n$; substitution of the two proved
absolute exponents and cancellation yields
$d(F_{n+1}/F_n)=p^{n+1}v_K(p)+p-1$.
Finally $p$ is an actual uniformizer of $\mathbb Q_p$.
Normalized valuation scaling along its original integer-ring map
gives $v_K(p)=e(K/\mathbb Q_p)$, proving the last assertion.
\end{proof}

\begin{lemma}[The relative different in the radical tower]
\label{ft:direct-radical-step}
\uses{ft:radical-layer-ramification,fw:prime-degree-nonnormal-index-profile}
\lean{FalseTateExtension.radicalLayer_step_own_ramificationIdx, FalseTateExtension.radicalLayer_step_own_different_order, FalseTateExtension.radicalLayer_step_apfIndex, FalseTateExtension.radicalLayer_step_apfPsi, FalseTateExtension.radicalLayer_step_apfRamificationBound}
\leanok
Let $p$ be any prime, let $K/\mathbb Q_p$ be finite, and put
$e_K=e(K/\mathbb Q_p)$. For compatible roots
$\pi_{n+1}^p=\pi_n$ of a uniformizer $\pi_0$ of $K$, set
$F_n=K(\pi_n)$. The field $F_n/K$ is totally ramified of degree
$p^n$, and $\pi_n$ is a uniformizer. The actual relative extension
$F_{n+1}/F_n$ has degree $p$ and
\[
 d(F_{n+1}/F_n)=p e_Kp^n+p-1,\qquad
 i(F_{n+1}/F_n)=t_n=\frac{e_Kp^{n+1}}{p-1}.
\]
Its intrinsic upper index is one through $t_n$ and $p$ after $t_n$.
No normality of this relative extension is required.
\end{lemma}
\begin{proof}
The polynomial $X^{p^n}-\pi_0$ is Eisenstein over the actual
integer ring of $K$. It is irreducible, its root is a uniformizer,
and its degree and ramification index are $p^n$. Alternatively,
its root has value $1/p^n$ in the normalization extending $v_K$;
the ramification index is therefore divisible by $p^n$, while
the polynomial bounds the degree by $p^n$. The fundamental degree
identity forces equality and residue degree one. Tower degrees
give $[F_{n+1}:F_n]=p$ and $e(F_n/\mathbb Q_p)=e_Kp^n$.
The uniformizer integral power basis makes $X^p-\pi_n$ the actual
monogenic polynomial for the relative extension. Its derivative
at $\pi_{n+1}$ is $p\pi_{n+1}^{p-1}$, whose normalized valuation
is $p e_Kp^n+p-1$. Equivalently the absolute different exponent is
$d_n=p^n v_K(p^n)+p^n-1$. Transitivity of the different gives
$d_{n+1}=d(F_{n+1}/F_n)+p d_n$; subtracting the two explicit
absolute exponents yields $p^{n+1}v_K(p)+p-1$, the same relative
exponent. The nonnormal prime-degree index-profile
identity $d=(p-1)(i+1)$ gives the displayed $t_n$ and profile.
To identify the intrinsic extended first bound explicitly, use any
actual separable-closure embedding of $F_{n+1}/F_n$.
Total ramification identifies its full APF-product index with its
intrinsic index for every nonnegative parameter. Since $p>1$,
that full product is the whole group exactly through $t_n$.
Thus $t_n$ belongs to the defining set of the extended supremum.
Every other parameter in that set is either negative, hence at
most $0\leq t_n$, or nonnegative, hence at most $t_n$ by the
proved index profile. The supremum is precisely $t_n$.
\end{proof}

\begin{lemma}[The index profile of the infinite radical tower]
\label{ft:direct-kummer-profile}
\uses{ft:direct-radical-step,fw:apf-nested-index-factor,fw:apf-finite-base-index,fw:apf-finite-base-herbrand,fw:apf}
\lean{FalseTateExtension.isAPF_kummerExtension, FalseTateExtension.isStrictlyAPF_kummerExtension, FalseTateExtension.kummerExtension_apfIndex_initial, FalseTateExtension.kummerExtension_apfIndex_interval, FalseTateExtension.kummerExtension_apfIndex_growth, FalseTateExtension.kummerExtension_apfPsi_break, FalseTateExtension.kummerExtension_apfPsi_div_index_ge, FalseTateExtension.apfSubgroup_kummerExtension_zero_eq_top, FalseTateExtension.kummerExtension_apfSubgroup_fieldRange_index}
\leanok
For the preceding tower set
$u_j=e_K(j+1/(p-1))$ for $j\geq1$. The extension
$K^{\mathrm{Kum}}/K=\bigcup_n F_n/K$ is APF. Its intrinsic index
$A(u)$ equals one on $[0,u_1]$ and equals $p^j$ on
$(u_j,u_{j+1}]$. In particular
\[
 A(u+e_K)\leq p A(u)\quad(u\geq0),\qquad
 \psi_{K^{\mathrm{Kum}}/K}(u_j)=\frac{e_Kp^j}{p-1}.
\]
This extension is strictly APF. For every actual embedding
$j:K^{\mathrm{Kum}}\hookrightarrow K^{\rm sep}$, its full
APF product at zero is $G_K$, and
$[G_K:G_{j(K^{\mathrm{Kum}})}G_K^u]=A(u)$ for every $u\geq0$.
\end{lemma}
\begin{proof}
Let $A_n$ and $\psi_n$ be the actual index and inverse Herbrand
function of $F_n/K$. The nested-index formula says
\[
 A_{n+1}(u)=A_n(u)\,
  a_{F_{n+1}/F_n}(\psi_n(u)).
\]
The last factor changes from one to $p$ when $\psi_n(u)$ passes
$t_n$. For $n=0$ the inverse function is the identity, so the
first new threshold is $u_1=t_0=e_Kp/(p-1)$.
Suppose the first $n$ thresholds are the displayed $u_j$.
Then $A_n=p^n$ after $u_n$, and direct integration of its
previous pieces gives $\psi_n(u_n)=t_{n-1}$.
Since $t_n-t_{n-1}=e_Kp^n$, the unique solution of
$\psi_n(u)=t_n$ is $u=u_n+e_K=u_{n+1}$.
The nested-index formula proves the next index profile, including
the endpoint because the relative prime-degree profile includes
its right endpoint. This completes the induction.

For each fixed $u$, the finite indices thus stabilize once
$u_n\geq u$. In the separable closure, the products of the
fixing groups with the fixed absolute upper group form a nested
sequence of closed subgroups. The bounded finite indices force
eventual equality. Their intersection is the product for the
union field: if an element is in every product, the corresponding
closed sets of possible factors in the compact absolute upper
group are nested and nonempty, so have nonempty intersection.
The limiting product is therefore open with the stabilized index.
This proves APF and the claimed infinite profile.
At zero every actual finite radical layer is totally ramified,
so its full APF product is $G_K$. The intersection identity
therefore gives the same equality for their union, without
requiring positivity of $v_K(p)$ for this zero-level assertion.
At every nonnegative parameter, stabilization identifies the
full infinite-product index with that of an actual finite
radical layer. Total ramification identifies the latter with
its intrinsic inertia index, already shown to equal $A(u)$.
Conjugacy of separable-closure embeddings transports both
statements to the specified $j$.

The thresholds are spaced by $e_K$, with the initial interval
at least that long; shifting by $e_K$ crosses at most one jump.
This proves the growth inequality. Integrating the profile gives
$\psi(u_1)=u_1=e_Kp/(p-1)$ and increments
$\psi(u_{j+1})-\psi(u_j)=e_Kp^j$, proving the formula at every
threshold. On $(u_j,u_{j+1}]$ the ratio $\psi(u)/A(u)$ is at
least $e_K/(p-1)$. At a threshold the index is smaller, so the
same bound holds. Its lower limit at infinity is positive,
which is strict APF by definition.
\end{proof}

\begin{lemma}[Actual finite cyclotomic displacements at every prime]
\label{ft:finite-cyclotomic-displacement}
\source{serre-1979-local-fields:page-0085}
\dcref{Serre, Chapter IV, Section 4, Propositions 17--18}
\uses{ft:cyclotomic-total-ramification,ft:finite-cyclotomic-character,fw:totally-ramified-uniformizer-generation,fw:generator-evaluation,fw:norm-valuation}
\lean{FalseTateExtension.norm_primitive_cyclotomic_sub_one, FalseTateExtension.normalizedValuation_primitive_cyclotomic_sub_one, FalseTateExtension.normalizedValuation_primitive_cyclotomic_pow_sub_one, FalseTateExtension.cyclotomic_exponent_decomposition, FalseTateExtension.normalizedValuation_primitive_cyclotomic_root, FalseTateExtension.normalizedValuation_cyclotomic_displacement_of_exponent, FalseTateExtension.normalizedValuation_cyclotomic_displacement, FalseTateExtension.cyclotomicUniformizer_adjoin_spectral, FalseTateExtension.ramificationOrder_cyclotomic_of_exponent, FalseTateExtension.ramificationOrder_cyclotomic}
\leanok
For every prime $p$, let $m\geq1$, $E_m=\mathbb Q_p(\zeta_m)$
and normalize $v_{E_m}$ to have value group $\mathbb Z$.
For every primitive $p^m$-th root $\zeta\in E_m$,
\[
 N_{E_m/\mathbb Q_p}(\zeta-1)=(-1)^{\varphi(p^m)}p,
 \qquad v_{E_m}(\zeta^{p^s}-1)=p^s\quad(0\leq s<m).
\]
The actual integer ring is generated by $\zeta_m-1$. If a
nonidentity automorphism has character $a\in(\mathbb Z/p^m)^\times$,
then its actual ramification order is
\[
 j(g)=p^{v_p(\widetilde a-1)},\qquad v_p(\widetilde a-1)<m,
\]
where $\widetilde a\in\{1,\ldots,p^m-1\}$ is the canonical
unit representative. This includes $p=2,m=1$, whose group is
trivial and whose signed norm is $-2$.
\end{lemma}
\begin{proof}
The minimal polynomial of $\zeta-1$ is
$\Phi_{p^m}(1+X)$, with constant term $p$ and degree
$\varphi(p^m)$. Its constant coefficient gives the signed norm.
The normalized norm identity and residue degree one show that
$\zeta-1$ is a uniformizer. For $s<m$, the root $\zeta^{p^s}$
is primitive of order $p^{m-s}$. The field it generates is a
totally ramified subfield of degree $\varphi(p^{m-s})$.
Restriction of normalized valuation therefore multiplies its
uniformizer order by
$\varphi(p^m)/\varphi(p^{m-s})=p^s$, proving the power formula.
The uniformizer generator theorem gives generation of the
actual integer ring.

For a nonidentity unit $a$, its representative differs from one.
Write $\widetilde a=1+p^s b$ with $p\nmid b$; its size implies
$s<m$. The element $\zeta_m^b$ is still primitive of order
$p^m$. The identity
\[
 \zeta_m^{\widetilde a}-\zeta_m
 =\zeta_m\bigl((\zeta_m^b)^{p^s}-1\bigr)
\]
and unit valuation of $\zeta_m$ give displacement order $p^s$.
The actual integral generator computes the ramification order
from precisely this displacement. The identity automorphism has
infinite order and is treated separately; no $p$-adic order of
zero is used.
\end{proof}

\begin{proposition}[The finite rational cyclotomic filtration at every prime]
\label{ft:finite-rational-cyclotomic-filtration}
\source{serre-1979-local-fields:page-0085}
\dcref{Serre, Chapter IV, Section 4, Proposition 18 and upper numbering}
\uses{ft:finite-cyclotomic-displacement,ft:finite-cyclotomic-character,fw:herbrand-functions,fw:real-lower-groups,fw:finite-upper-groups}
\lean{FalseTateExtension.mem_lowerRamificationGroup_cyclotomic_iff, FalseTateExtension.mem_lowerRamificationGroupReal_cyclotomic_iff, FalseTateExtension.cyclotomicUnitReduction_ker_index, FalseTateExtension.cyclotomicUnitReduction_ker_card, FalseTateExtension.mem_cyclotomicUnitReduction_ker_iff, FalseTateExtension.lowerRamificationGroup_cyclotomic_eq_comap_unitReduction, FalseTateExtension.lowerRamificationGroup_cyclotomic_index, FalseTateExtension.lowerRamificationGroup_cyclotomic_zero, FalseTateExtension.lowerRamificationGroup_cyclotomic_map_character, FalseTateExtension.lowerRamificationGroup_cyclotomic_card, FalseTateExtension.lowerRamificationGroupReal_cyclotomic_threshold, FalseTateExtension.lowerRamificationGroupReal_cyclotomic_constant, FalseTateExtension.lowerRamificationGroupReal_cyclotomic_card_on, FalseTateExtension.integral_cyclotomic_lower_slope, FalseTateExtension.herbrandFunction_cyclotomic_threshold, FalseTateExtension.upperRamificationGroup_cyclotomic_eq_comap_unitReduction, FalseTateExtension.upperRamificationGroup_cyclotomic_index, FalseTateExtension.upperRamificationGroup_cyclotomic_card, FalseTateExtension.upperRamificationGroup_cyclotomic_zero, FalseTateExtension.upperRamificationGroup_cyclotomic_eq_bot}
\leanok
In $E_m/\mathbb Q_p$ let $G=\operatorname{Gal}(E_m/\mathbb Q_p)$
and identify it by its actual character. Then $G_0=G$, and for
$0\leq r\leq m$,
\[
 G_{p^r-1}=\ker\bigl(G\longrightarrow(\mathbb Z/p^r)^\times\bigr),
 \quad [G:G_{p^r-1}]=\varphi(p^r),\quad
 \varphi_{E_m/\mathbb Q_p}(p^r-1)=r.
\]
For $r\geq1$ the displayed subgroup has order $p^{m-r}$.
The real lower group is constant on
$(p^r-1,p^{r+1}-1]$. For $0\leq r<m$ and $r<u\leq r+1$,
\[
 G^u=\ker\bigl(G\longrightarrow(\mathbb Z/p^{r+1})^\times\bigr),
 \quad [G:G^u]=\varphi(p^{r+1}),\quad |G^u|=p^{m-r-1}.
\]
Moreover $G^0=G$ and $G^u=1$ for $u>m-1$.
All assertions hold at every prime and retain their closed right
endpoints. Here $m$ is the positive cyclotomic level.
\end{proposition}
\begin{proof}
The displacement formula gives $g\in G_t$ exactly when $g=1$
or $t+1\leq p^{v_p(\widetilde a-1)}$. At $t=p^r-1$ this is
exactly $a\equiv1\pmod{p^r}$, including $r=0$ and the identity.
Reduction of units modulo prime powers is surjective: every unit
class has a representative prime to $p$, which remains a unit
at the higher level. Thus its kernel has index $\varphi(p^r)$;
for $r>0$ division of the two group orders gives $p^{m-r}$.
The same displacement powers show constancy on each indicated
real interval, with the closed right endpoint.

On $(p^r-1,p^{r+1}-1]$, for $r<m$, the lower subgroup has order
$p^{m-r-1}$. The interval length is $(p-1)p^r$ and $|G_0|$ is
$(p-1)p^{m-1}$, so the integral of $|G_t|/|G_0|$ over it is
exactly one. Starting at $\varphi(0)=0$ and adding consecutive
integrals gives $\varphi(p^r-1)=r$ for $0\leq r\leq m$.
Strict increase and the actual inverse identities then put
$\psi(u)$ in $(p^r-1,p^{r+1}-1]$ exactly when $r<u\leq r+1$.
The upper-group definition gives the claimed kernels and orders.
For $u>m-1$ the group is trivial through $u=m$, and monotonicity
makes it trivial thereafter. The equality at zero follows from
full inertia. These arguments include the trivial first dyadic
layer without omitting any interval.
\end{proof}

\begin{lemma}[Cyclotomic upper-index growth over a finite base]
\label{ft:direct-cyclotomic-growth}
\uses{ft:cyclotomic-total-ramification,ft:finite-cyclotomic-character,ft:finite-rational-cyclotomic-filtration,ft:cyclotomic-image-open,fw:absolute-finite-base-upper,fw:apf-finite-extension,fw:apf-finite-base,fw:apf-finite-slope,fw:strict-apf-finite-base,fw:inertia-relative-index,fw:abstract-normed-finite-base,fw:finite-base-eventual-ramification-slope,fw:apf-finite-stages}
\lean{FalseTateExtension.mem_upperRamificationGroupInfinite_cyclotomic_iff, FalseTateExtension.upperRamificationGroupInfinite_cyclotomic_eq_ker, FalseTateExtension.upperRamificationGroupInfinite_cyclotomic_zero, FalseTateExtension.upperRamificationGroupInfinite_cyclotomic_of_nonpos, FalseTateExtension.upperRamificationGroupInfinite_cyclotomic_eq_character_ker, FalseTateExtension.upperRamificationGroupInfinite_cyclotomic_map_character, FalseTateExtension.upperRamificationGroupInfinite_cyclotomic_index, FalseTateExtension.upperRamificationGroupInfinite_cyclotomic_finiteIndex, FalseTateExtension.isAPF_cyclotomicExtension_padic, FalseTateExtension.apfIndex_cyclotomicExtension_padic, FalseTateExtension.upperRamificationGroupInfinite_cyclotomic_index_add_one, FalseTateExtension.apfIndex_cyclotomicExtension_padic_add_one, FalseTateExtension.isStrictlyAPF_cyclotomicExtension_padic, FalseTateExtension.mem_cyclotomicLayer_of_pow_eq_one, FalseTateExtension.rationalCyclotomicEmbedding, FalseTateExtension.cyclotomicRootOfUnity_mem_rationalCyclotomicEmbedding, FalseTateExtension.rationalCyclotomicEmbedding_adjoin_base, FalseTateExtension.rationalCyclotomicEmbedding_finiteDimensional, FalseTateExtension.isStrictlyAPF_cyclotomicExtension_over_padic, FalseTateExtension.isStrictlyAPF_cyclotomicExtension, FalseTateExtension.isAPF_cyclotomicExtension, FalseTateExtension.cyclotomicCharacter_map_upper_autCongr, FalseTateExtension.cyclotomicExtension_map_upper_from_algebraicClosure, FalseTateExtension.cyclotomicCharacter_map_upperRamificationGroupInfinite, FalseTateExtension.cyclotomicCharacter_map_upper_finiteBase, FalseTateExtension.cyclotomicExtension_map_upper_finiteBase, FalseTateExtension.exists_cyclotomicExtension_upper_principal_finiteBase, FalseTateExtension.exists_cyclotomicExtension_upper_index_finiteBase, FalseTateExtension.exists_upperRamificationGroupInfinite_cyclotomic_index_add_ramificationIdx}
\leanok
For every prime $p$ and finite $K/\mathbb Q_p$, the extension
$C=K^{\mathrm{cyc}}/K$ is strictly APF. Write
$H=\Gal(C/K)$ in its actual cyclotomic character and let
$B(u)=[H:H^u]$ for $u\geq0$. These are finite indices and,
for all sufficiently large $u$,
\[
 B(u+e_K)=p B(u).
\]
\end{lemma}
\begin{proof}
First take $K=\mathbb Q_p$. In the layer generated by a primitive
$p^n$-th root $\zeta$, the element $\zeta-1$ is a uniformizer.
For a nonidentity exponent $a$, put $j=v_p(a-1)<n$.
The primitive-root valuation gives
$v(\zeta^{a-1}-1)=p^j$, hence its lower break is $p^j-1$.
For odd $p$, the first positive lower interval has length $p-1$
and inertia index $p-1$, and subsequent intervals have length
$(p-1)p^j$ and index $(p-1)p^j$. Every such interval has upper
length one. For $p=2$, all units already lie in the first
principal-unit group: the initial positive interval has length
one and index one, and each subsequent interval again has upper
length one. The level $n=1$ at $p=2$ is trivial. Thus in the
infinite rational tower the actual character image at $u>0$ is
$1+p^{\lceil u\rceil}\mathbb Z_p$ at every prime. These open
groups prove APF. Their index grows by a factor $p$ on shifting
$u$ by one. Integrating over the preceding unit interval gives
a positive eventual lower bound for $\psi(u)/a(u)$; explicitly
$\psi(u)\geq a(u-1)\geq a(u)/p$ for $u$ sufficiently large.
This proves rational strict APF.

Put $C_0=\mathbb Q_p(\mu_{p^\infty})$ and $N=KC_0$ inside the
same algebraic closure. The extension $N/C_0$ is finite, so finite-top
strictness gives strict APF for $N/\mathbb Q_p$; finite-base
strictness then gives it for $N/K=C/K$. This handles a base $K$
which is not contained in $C_0$. To obtain the
quantitative assertion, put $v=\varphi_{K/\mathbb Q_p}(u)$.
After the last break of a finite Galois closure of $K/\mathbb Q_p$,
the rational absolute upper group is contained in the actual fixing
group of $K$. Absolute finite base change therefore identifies its
image in the cyclotomic character with the actual group $H^u$,
where $u=\psi_{K/\mathbb Q_p}(v)$. The finite-base inverse function
has eventual slope $e_K$: once the absolute upper group lies in
the fixing group of $K$, its intrinsic index is the fixed index
of the base inertia modulo its intersection with that fixing
group. Projection to a finite Galois closure and the actual
inertia-relative index formula identify this index with
$e(K/\mathbb Q_p)=e_K$. Integrating this constant tail gives
$\psi_{K/\mathbb Q_p}(v)=e_Kv+c$ for a fixed $c$. Replacing $u$
by $u+e_K$ consequently replaces $v$ by $v+1$. The rational
principal-unit formula gives precisely the factor $p$ in the
displayed full index $[H:H^u]$. More precisely, put
$D=[\mathbb Z_p^\times:\chi(H)]$, which is a positive finite
integer by openness of the actual character image. Equality of the
upper character images gives
$B_K(\psi_{K/\mathbb Q_p}(v))D=B_{\mathbb Q_p}(v)$.
The two formulas at $v$ and $v+1$, together with the rational
shift identity, give equal products with $D$; cancellation of
this positive integer proves the assertion.
\end{proof}

\begin{lemma}[The actual false Tate affine indices]
\label{ft:direct-affine-index-comparison}
\source{wintenberger:page-0005}
\dcref{Local application of the closed affine index comparison to the actual upper groups}
\uses{ft:affine-embedding,ft:cyclotomic-image-open,ft:closed-affine-index-comparison,ft:affine-zero-kernel-isomorphism,ft:direct-kummer-profile,ft:direct-cyclotomic-growth,fw:infinite-upper-restriction,fw:galois-apf-criterion}
\lean{FalseTateExtension.falseTateAffineImage, FalseTateExtension.falseTateAffineEquiv, FalseTateExtension.falseTateAffineEquiv_apply, FalseTateExtension.falseTateAffineEquiv_projection, FalseTateExtension.falseTateAffineImage_projection_range, FalseTateExtension.kummerExtension_le_falseTateExtension, FalseTateExtension.falseTateKummerInclusion, FalseTateExtension.falseTateKummerFixing, FalseTateExtension.mem_falseTateKummerFixing_iff, FalseTateExtension.falseTateAffineEquiv_mem_pureUnits, FalseTateExtension.falseTateKummerFixing_map_affine, FalseTateExtension.falseTateAffineImage_projection_infinite, FalseTateExtension.padicCyclotomicCharacter_injective_of_translationKernel_eq_bot, FalseTateExtension.falseTate_cyclotomicFixing_eq_bot_of_translationKernel_eq_bot, FalseTateExtension.falseTateExtension_eq_cyclotomicExtension_of_translationKernel_eq_bot, FalseTateExtension.apfSubgroup_fieldRange_index_eq_falseTateFixing_sup_upper, FalseTateExtension.falseTateCharacter_map_upperRamificationGroupInfinite, FalseTateExtension.cyclotomicRestriction_map_upperRamificationGroupInfinite, FalseTateExtension.falseTateAffine_upper_projection_index, FalseTateExtension.falseTateAffine_upper_radical_index, FalseTateExtension.exists_falseTate_upper_index_comparison, FalseTateExtension.falseTate_upper_index_of_translationKernel_eq_bot}
\leanok
Let $G\subseteq\mathbb Z_p\rtimes\mathbb Z_p^\times$ be the
actual closed affine image of the false Tate group, let $H$ be
its unit projection, $M$ its translation kernel, and
$S=G\cap(\{0\}\rtimes\mathbb Z_p^\times)$.
Write $U(u)$ for its actual upper ramification group, and put
\[
 A(u)=a_{K^{\mathrm{Kum}}/K}(u),\qquad
 B(u)=[\operatorname{Gal}(K^{\mathrm{cyc}}/K):
                 \operatorname{Gal}(K^{\mathrm{cyc}}/K)^u].
\]
Then $S$ is the actual radical fixing subgroup,
$[G:SU(u)]=A(u)$ for $u\geq0$, and
$[H:\chi(U(u))]=B(u)$ for $u\geq-1$.
If $M\ne0$, there are positive integers $C,D$ such that, for
all $u\geq0$, $[G:U(u)]$ is finite and
\[
 A(u)B(u)\leq D[G:U(u)],\qquad
 [G:U(u)]\leq C A(u)B(u).
\]
If $M=0$, the actual false Tate extension equals its
cyclotomic subextension and $[G:U(u)]=B(u)$ for $u\geq-1$.
\end{lemma}
\begin{proof}
The faithful continuous affine map from the original compact Galois
group gives a topological group isomorphism onto its closed image.
An automorphism fixes every compatible radical exactly when each
finite Kummer exponent vanishes, equivalently when its $p$-adic
translation coordinate is zero. Thus its actual radical fixer
maps exactly to $S$.
Ordinary upper restriction from a common algebraic closure identifies
the upper character images for the actual false Tate and cyclotomic
fields. The character on the cyclotomic group is injective. Its image
has positive finite index in the full unit group, so cancellation
of that common positive index gives $[H:\chi(U(u))]=B(u)$.
Likewise, restriction to the original false Tate field identifies the
full radical APF product index with $[G:SU(u)]$. The radical tower
is totally ramified, so its full product index is exactly the
intrinsic inertia index $A(u)$.
The actual unit projection is infinite: the full $p$-adic unit
group is infinite, while the actual cyclotomic image has finite
index. For $M\ne0$, both $A(u)$ and $B(u)$ are positive finite
by the already proved radical and cyclotomic APF assertions.
Apply the closed affine index comparison to the actual closed
normal group $U(u)$. It produces finite index and the two bounds.
For $M=0$, the actual cyclotomic character is injective, hence
its original field restriction has trivial fixing subgroup.
Infinite Galois correspondence makes the cyclotomic subfield the
entire original false Tate field. The projection isomorphism and
the already proved projection-index equality also give the final
index formula.
\end{proof}

\begin{lemma}[Bounded index growth implies strict APF]
\label{ft:direct-index-growth-strictness}
\uses{fw:apf-bounded-index-growth}
\lean{FontaineWintenberger.isStrictlyAPF_of_apfIndex_boundedGrowth}
\leanok
Let $L/K$ be APF, with intrinsic inertia index $a(u)$.
If there are $h>0$, $Q\geq1$ and $U\geq0$ such that
$a(u+h)\leq Q a(u)$ for every $u\geq U$, then $L/K$ is
strictly APF.
\end{lemma}
\begin{proof}
For $v\geq U+h$, monotonicity of the actual index and its integral
formula give
\[
 \psi_{L/K}(v)=\int_0^v a(t)\,dt
 \geq\int_{v-h}^v a(t)\,dt
 \geq h a(v-h)\geq\frac hQ a(v).
\]
The first inequality uses nonnegativity of the integral over
$[0,v-h]$. Dividing by the positive finite index shows that
the lower limit of the strictness ratio is at least $h/Q>0$.
This proves strict APF directly from its definition, including
extensions with a finite initial unramified or tame part.
\end{proof}

\begin{theorem}[Strict arithmetic profiniteness]
\label{ft:strict-apf}
\uses{ft:direct-kummer-profile,ft:direct-cyclotomic-growth,ft:direct-affine-index-comparison,ft:direct-index-growth-strictness,fw:galois-apf-criterion,fw:infinite-upper-restriction,fw:apf-inertia-indices}
\lean{FalseTateExtension.isAPF_falseTateExtension, FalseTateExtension.exists_falseTate_upper_index_boundedGrowth, FalseTateExtension.falseTate_upper_index_eq_apfIndex_mul_inertia_index, FalseTateExtension.exists_falseTate_apfIndex_boundedGrowth, FalseTateExtension.isStrictlyAPF_falseTateExtension, FalseTateExtension.isStrictlyAPF_falseTateExtension_of_finiteDimensional, FalseTateExtension.isAPF_falseTateExtension_of_finiteDimensional, FalseTateExtension.isStrictlyAPF_kummerExtension, FalseTateExtension.isStrictlyAPF_cyclotomicExtension}
\leanok
\source{wintenberger:page-0005}
\dcref{Wintenberger, Definition 1.2.1 and Proposition 1.2.3; direct radical and affine argument}
For every prime $p$, the extensions $K^{\mathrm{FT}}_\infty/K$,
$K^{\mathrm{cyc}}/K$, and $K^{\mathrm{Kum}}/K$ are strictly APF.
\end{theorem}
\begin{proof}
The radical and cyclotomic assertions follow from the direct profiles
just proved. Write $G$ for the actual false Tate group in its closed
affine image. If its translation kernel is zero, the extension
equals its cyclotomic subextension and the assertion follows.
Otherwise use the fixed constants $C,D$ of the affine index
comparison. For its actual closed normal upper group $I(u)$,
ordinary upper restriction identifies
\[
 [G:S I(u)]=A(u),\qquad [H:\chi(I(u))]=B(u),
\]
with the actual radical and cyclotomic full indices. The former
equals the intrinsic radical index because that tower is totally
ramified. Both are finite by the two APF assertions. The comparison
therefore proves $[G:I(u)]$ finite for every $u\geq0$; it is
also finite for $-1\leq u<0$ by monotonicity. The intrinsic Galois
criterion gives APF of the actual false Tate extension.

For all sufficiently large $u$, the two explicit growth bounds give
\[
 [G:I(u+e_K)]\leq C A(u+e_K)B(u+e_K)
 \leq Cp^2 A(u)B(u)\leq CDp^2[G:I(u)].
\]
The actual intrinsic inertia index is
$a(u)=[G:I(u)]/[G:I(0)]$. Its fixed positive denominator cancels,
so the same bound holds for $a$. Apply bounded-index-growth
strictness with $h=e_K>0$ and $Q=CDp^2$.
This argument uses the prime-degree different identity, the actual
radical profile and the affine action at every prime; it does not
require total ramification of the full tower over a general base.
\end{proof}

\subsection{Quotient compatibility in a totally ramified tower}

For the finite calculations we use the totally ramified specializations of
Lemmas~\ref{fw:totally-ramified-uniformizer-generation},
\ref{fw:integral-fraction-generator} and~\ref{fw:generator-evaluation},
the coset average of Lemma~\ref{fw:coset-ramification-average}, and
Theorems~\ref{fw:herbrand-lower-quotients},
\ref{fw:herbrand-transitivity} and~\ref{fw:herbrand-quotients}.
These results apply to the actual local fields here. In a totally ramified
tower inertia is the full Galois group, so the coset-average denominator is
the order of the relative Galois group. These are the totally ramified
cases of Serre\textquotesingle s quotient theorem
\cite[Chapter IV, Sections 1 and 3]{Serre1979}.

\section{Finite diagonal affine quotients}

\begin{lemma}[The actual diagonal finite affine quotient]
\label{ft:diagonal-finite-affine-structure}
\uses{ft:galois-structure,ft:finite-coordinates,ft:galois-layers,ft:extensions,fw:profinite-galois-foundations}
\lean{FalseTateExtension.finiteDiagonalGaloisEquiv, FalseTateExtension.finrank_falseTateLayer_diagonal, FalseTateExtension.finrank_falseTateLayer_over_cyclotomic, FalseTateExtension.finiteDiagonalGaloisEquiv_mem_translations, FalseTateExtension.finiteDiagonalGaloisEquiv_mem_units}
\leanok
\source{liu-2008-lattices:page-0022,liu-2008-lattices:page-0023}
\dcref{Finite-quotient consequence of Liu, Lemma 5.1.2}
Let $p\geq3$, let $K/\mathbb Q_p$ be finite, and let $\pi$ be a
uniformizer of $K$. For $n\geq1$ put
\[
 C_{K,n}=K(\zeta_n),\qquad D_{K,n}=K(\pi_n,\zeta_n),
\]
and let $H_{K,n}\subseteq(\mathbb Z/p^n\mathbb Z)^\times$ be the
actual image of the cyclotomic character of $\operatorname{Gal}(C_{K,n}/K)$.
The action on these two generators identifies the original finite
Galois group with
\[
 \operatorname{Gal}(D_{K,n}/K)
 \simeq (\mathbb Z/p^n\mathbb Z)\rtimes H_{K,n},
 \qquad (b,u)(c,v)=(b+uc,uv),
\]
where $(b,u)$ acts by
$\pi_n\mapsto\zeta_n^b\pi_n$ and $\zeta_n\mapsto\zeta_n^u$.
Consequently
\[
 [D_{K,n}:K]=p^n[C_{K,n}:K],\qquad
 [D_{K,n}:C_{K,n}]=p^n.
\]
In these same coordinates the subgroup fixing $C_{K,n}$ is the full
translation group and the subgroup fixing $K(\pi_n)$ is the pure-unit
subgroup $\{0\}\rtimes H_{K,n}$.
\end{lemma}
\begin{proof}
The field $D_{K,n}$ is the splitting field of $X^{p^n}-\pi$ over $K$:
it contains all its roots $\zeta_n^b\pi_n$, and conversely these roots
supply $\pi_n$ and their quotients supply $\zeta_n$.
The polynomial is separable in characteristic zero, so the extension is
finite Galois. Its normalized finite action coordinates are injective,
since fixing both generators fixes the whole field. Applying two
automorphisms to the generators gives the displayed affine product.

Write $H_K$ for the actual infinite cyclotomic character image.
Every automorphism of $C_{K,n}/K$ extends to the cyclotomic union,
by Galois restriction; hence the reduction of $H_K$ modulo $p^n$
is exactly $H_{K,n}$. Given $b\in\mathbb Z/p^n\mathbb Z$ and
$u\in H_{K,n}$, lift $b$ to $\beta\in\mathbb Z_p$ and lift $u$
to $h\in H_K$. The proved infinite Galois structure gives an actual
automorphism of the original false Tate union with coordinates
$(\beta,h)$. Its restriction to the stable diagonal field $D_{K,n}$
has exactly the prescribed finite coordinates $(b,u)$.
Thus the finite action map is onto the entire displayed semidirect
product, rather than merely onto a subgroup.

The finite cyclotomic character is injective, since $\zeta_n$
generates $C_{K,n}$. Thus $|H_{K,n}|=[C_{K,n}:K]$.
The cardinality of the finite affine group and the finite Galois degree
formula give the first degree equality; degree multiplicativity gives
the second. Finally primitivity of $\zeta_n$ and nonzeroness of
$\pi_n$ show that fixing $\pi_n$ is equivalent to $b=0$, and fixing
$\zeta_n$ is equivalent to $u=1$. These are the asserted original
fixing subgroups. This finite-quotient argument uses the full normalized
infinite coordinates and the actual finite cyclotomic image.
\end{proof}

\section{Finite radical layers over the rational p-adic field}

Throughout this section $p$ is odd and $a\in\mathbb Z$ satisfies
$v_p(a)=1$. The calculation follows the extensions considered in
\cite[Section 6]{Viviani2004}. We recover the translation depth of each
actual upper group from the intrinsic indices of the pure radical
subfields. The finite Galois-overfield index identity relates these two
quantities, including every right endpoint. This replaces the inductive
mixed-tower calculation while keeping the original fields and groups.

\begin{notation}[Finite radical fields and affine coordinates]
\label{ft:finite-radical-setup}
\uses{ft:extensions}\lean{FalseTateExtension.finiteTranslationSubgroup, FalseTateExtension.finiteTranslationSubgroup_card, FalseTateExtension.finiteUnitSubgroup, FalseTateExtension.finiteAffineSubgroup, FalseTateExtension.irreducible_integerParameter}
\leanok
Choose compatible elements $\alpha_s$ with $\alpha_s^{p^s}=a$ and
$\alpha_{s+1}^p=\alpha_s$, and compatible primitive roots $\zeta_n$ of
order $p^n$. For $n\geq1$ and $1\leq s\leq n$ put
\[
 F_s=\mathbb Q_p(\alpha_s),\qquad
 C_n=\mathbb Q_p(\zeta_n),\qquad
 L_n(a)=C_n(\alpha_n),\qquad
 G_n=\operatorname{Gal}(L_n(a)/\mathbb Q_p).
\]
Put $F_0=\mathbb Q_p$. In the affine group
$A_n=(\mathbb Z/p^n\mathbb Z)\rtimes(\mathbb Z/p^n\mathbb Z)^\times$
the product is $(b,u)(c,v)=(b+uc,uv)$. Write
\[
 T_k=p^k\mathbb Z/p^n\mathbb Z\quad(0\leq k\leq n),\qquad
 U_0=(\mathbb Z/p^n\mathbb Z)^\times,\qquad
 U_j=1+p^j\mathbb Z/p^n\mathbb Z\quad(j\geq1).
\]
Thus $T_n=\{0\}$ and $U_j=\{1\}$ for $j\geq n$.
The notation $C(p^b)$ denotes $T_{n-b}$, a cyclic group of order $p^b$.
\end{notation}

\begin{lemma}[Finite Galois structure]
\label{ft:finite-galois-structure}
\uses{ft:extensions,ft:diagonal-finite-affine-structure,ft:finite-cyclotomic-character}\lean{FalseTateExtension.finiteDiagonalGaloisEquiv_padic, FalseTateExtension.finiteDiagonalGaloisEquiv_padic_left, FalseTateExtension.finiteDiagonalGaloisEquiv_padic_right, FalseTateExtension.finrank_falseTateLayer_diagonal_padic, FalseTateExtension.finiteDiagonalGaloisEquiv_padic_mem_translations, FalseTateExtension.finiteDiagonalGaloisEquiv_padic_mem_units}
\leanok
\source{liu-2008-lattices:page-0022,liu-2008-lattices:page-0023}
\dcref{Rational-base specialization of the finite affine quotient}
For $n\geq1$, the extension $L_n(a)/\mathbb Q_p$ is Galois and the rules
\[
 \sigma(\alpha_n)=\zeta_n^b\alpha_n,\qquad
 \sigma(\zeta_n)=\zeta_n^u
\]
identify $G_n$ with $A_n$. In particular,
$[L_n(a):\mathbb Q_p]=p^n(p-1)p^{n-1}$.
The full translation subgroup is precisely
$\operatorname{Gal}(L_n(a)/C_n)$, and the pure-unit subgroup
$\{0\}\rtimes U_0$ is precisely $\operatorname{Gal}(L_n(a)/F_n)$.
\end{lemma}
\begin{proof}
The element $a$ has normalized valuation one in $\mathbb Q_p$, so it
is a uniformizer. Apply the preceding diagonal quotient lemma with
$K=\mathbb Q_p$ and $\pi=a$, using the specified compatible roots.
The finite rational cyclotomic character is onto the whole unit group,
so $H_{K,n}=U_0$. Its order is $(p-1)p^{n-1}$.
The affine isomorphism, its degree formula, and its actual pure-unit
fixing subgroup give every assertion, including at $n=1$.
\end{proof}

\begin{lemma}[Commutators act trivially on a rational residue extension]
\label{ft:finite-rational-residue-commutators}
\uses{fw:local-inertia-size,fw:local-dvr-degree-identity}
\lean{FalseTateExtension.finite_spectralResidue_padic, FalseTateExtension.isCyclic_rationalResidueGaloisImage, FalseTateExtension.commutator_mem_inertia_padic}
\leanok
\source{serre-1979-local-fields:page-0069}
\dcref{Finite-field consequence of the residue action in Chapter IV, Section 1}
For a finite Galois extension $E/\mathbb Q_p$, let $G$ be its actual
Galois group and let $I=G_0$ be its actual inertia subgroup.
The image of the action of $G$ on $k_E$ is cyclic, and every
commutator of $G$ belongs to $I$.
\end{lemma}
\begin{proof}
The residue field $k_E$ is a finite extension of $k_{\mathbb Q_p}=\mathbb F_p$,
since its degree is finite by the local degree identity. Write
$|k_E|=p^f$. The Frobenius $x\mapsto x^p$ is an automorphism,
and its order is exactly $f$: its $f$th power is the identity,
whereas for $0<r<f$ the polynomial $X^{p^r}-X$ has fewer than
$p^f$ roots and hence cannot vanish on all of $k_E$.
The finite field extension is simple of degree $f$, so an automorphism
is determined by a root of the minimal polynomial of a generator.
There are at most $f$ such automorphisms. The $f$ distinct Frobenius
powers therefore give its entire automorphism group, which is cyclic.

Every automorphism of $E/\mathbb Q_p$ preserves the unique extended
valuation and hence induces an automorphism of $k_E$ fixing $\mathbb F_p$.
Its kernel consists exactly of those automorphisms satisfying
$g(x)-x\in\mathfrak m_E$ for every $x\in\mathcal O_E$; this is the
original ideal definition of $G_0$.
The residue image is a subgroup of the cyclic group just described,
so is cyclic and in particular abelian. Every commutator is consequently
in the kernel $I$. No surjectivity of the residue action is needed.
\end{proof}

\begin{theorem}[Total ramification]
\label{ft:total-ramification}
\uses{ft:finite-galois-structure,ft:cyclotomic-total-ramification,ft:finite-rational-residue-commutators,fw:inertia-residue-index}
\lean{FalseTateExtension.lowerRamificationGroup_falseTateLayer_zero, FalseTateExtension.falseTateLayer_spectral_residueDegree, FalseTateExtension.falseTateLayer_spectral_ramificationIdx, FalseTateExtension.falseTateExtension_residue_algebraMap_surjective, FalseTateExtension.falseTateExtension_residueDegree_padic, FalseTateExtension.finiteSubextension_falseTate_residueDegree_padic, FalseTateExtension.finiteSubextension_falseTate_ramificationIdx_padic}
\leanok\source{viviani-2004-radical-extensions:page-0032}
\dcref{Viviani, Corollary 6.2; odd-prime affine proof}
For every odd prime $p$ and every $n\geq1$, $L_n(a)/\mathbb Q_p$ is
totally ramified. Consequently $\bigcup_nL_n(a)/\mathbb Q_p$ is totally
ramified, including the false Tate extension obtained by taking $a=p$.
\end{theorem}
\begin{proof}
Let $I=(G_n)_0$ be the original inertia subgroup. Every commutator
belongs to $I$ by the preceding finite residue-action lemma.
Since $p$ is odd, $2$ is a unit modulo $p^n$, and the actual affine
coordinates give
\[
 [(0,2),(b,1)]=(b,1)\qquad(b\in T_0).
\]
Thus $T_0\subseteq I$. Its actual fixed field is $C_n$, by the
finite affine fixing-subgroup identity.
The residue-index formula in the actual finite Galois overfield
$L_n(a)$ gives
\[
 1=f(C_n/\mathbb Q_p)=[G_n:T_0I]=[G_n:I],
\]
where the first equality is cyclotomic total ramification.
Hence $I=G_n$. The same residue-index formula for $L_n(a)$ itself
now gives $f(L_n(a)/\mathbb Q_p)=1$, proving total ramification.
Every finite subextension of the directed diagonal union is contained
in some $L_n(a)$: put a finite set of its field generators in one
common diagonal layer. Its residue field is therefore also $\mathbb F_p$.
This proves the infinite assertion with the same explicit odd-prime
restriction.
\end{proof}

\subsection{Intrinsic indices of the pure radical fields}

\begin{lemma}[The complete finite pure-radical index profile]
\label{ft:finite-radical-index-profile}
\uses{ft:finite-radical-setup,ft:direct-radical-step,fw:apf-finite-base-index,fw:apf-finite-base-herbrand}
\lean{FalseTateExtension.radicalTranslationDepth, FalseTateExtension.radicalLayer_apfIndex_eq_pow_translationDepth, FalseTateExtension.radicalTranslationDepth_interval, FalseTateExtension.radicalTranslationDepth_initial, FalseTateExtension.radicalTranslationDepth_tail}
\leanok
\source{viviani-2004-radical-extensions:page-0033,viviani-2004-radical-extensions:page-0034}
\dcref{Intrinsic-index adaptation of the prime-radical tower calculation}
Put $b_r=r+1/(p-1)$ for $r\geq1$. For every $n\geq0$ and $u\geq0$,
the intrinsic index of the original pure radical field is
\[
 a_{F_n/\mathbb Q_p}(u)=p^{m_n(u)},\qquad
 m_n(u)=\#\{r:1\leq r\leq n,\ b_r<u\}.
\]
Equivalently it is one on $[0,b_1]$, is $p^{\min(n,j)}$ on
$(b_j,b_{j+1}]$ for $j\geq1$, and is $p^n$ after $b_n$ when $n\geq1$.
The assertion for $n=0$ is the constant index one.
\end{lemma}
\begin{proof}
The relative extension $F_{n+1}/F_n$ is a totally ramified degree-$p$
extension with intrinsic index one through
$t_n=p^{n+1}/(p-1)$ and $p$ after $t_n$, by the direct radical-step lemma
with $e(\mathbb Q_p/\mathbb Q_p)=1$. This statement concerns the original
nonnormal pure radical fields; no Galois hypothesis on $F_n/\mathbb Q_p$
is used.

Write $A_n=a_{F_n/\mathbb Q_p}$ and
$\Psi_n=\psi_{F_n/\mathbb Q_p}$. The finite-base index identity gives
\[
 A_{n+1}(u)=A_n(u)
 \begin{cases}
 1,&\Psi_n(u)\leq t_n,\\
 p,&\Psi_n(u)>t_n.
 \end{cases}
\]
For $n=0$, $A_0=1$ and $\Psi_0$ is the identity, so the first
threshold is $b_1=p/(p-1)$.
Suppose the displayed profile has been proved for $n$.
Its integral at the next proposed threshold is
\[
 \Psi_n(b_{n+1})
 =b_1+\sum_{r=1}^{n}p^r(b_{r+1}-b_r)
 =\frac p{p-1}+\sum_{r=1}^{n}p^r
 =\frac{p^{n+1}}{p-1}=t_n.
\]
The empty sum gives the same equality when $n=0$.
The intrinsic inverse Herbrand function is strictly increasing.
Thus the new factor is one exactly for $u\leq b_{n+1}$ and is $p$
exactly for $u>b_{n+1}$. This proves the formula for $n+1$ by induction.
Counting the strictly smaller thresholds gives the equivalent interval
formulas, with the asserted finite-level truncation and right endpoints.
The thresholds are the same as in Viviani's mixed Galois calculation;
here they are obtained by the intrinsic nonnormal finite-base identities.
\end{proof}

\subsection{The upper filtration}

\begin{lemma}[The cyclotomic upper filtration]
\label{ft:cyclotomic-upper-filtration}
\uses{ft:finite-radical-setup,ft:finite-rational-cyclotomic-filtration}
\lean{FalseTateExtension.upperRamificationGroup_cyclotomic_zero, FalseTateExtension.upperRamificationGroup_cyclotomic_eq_comap_unitReduction, FalseTateExtension.upperRamificationGroup_cyclotomic_eq_bot}
\leanok\source{serre-1979-local-fields:page-0085}
\dcref{Serre, Chapter IV, Section 4}
The upper filtration on $\operatorname{Gal}(C_n/\mathbb Q_p)=U_0$ is
$U_0$ at $u=0$ and $U_{\lceil u\rceil}$ for $u>0$.
\end{lemma}
\begin{proof}
The actual finite cyclotomic character is the identification with $U_0$.
The proved rational cyclotomic filtration gives the whole group at zero
and, for every positive real parameter $u$, exactly the units congruent
to one modulo $p^{\lceil u\rceil}$ under this same character.
This is $U_{\lceil u\rceil}$ by definition, with trivial group when
$\lceil u\rceil\geq n$. In particular the first layer has trivial
positive upper groups, and every right endpoint is included as stated.
\end{proof}

\begin{lemma}[Normal subgroups in affine coordinates]
\label{ft:normal-affine-subgroup}\uses{ft:finite-radical-setup}\lean{FalseTateExtension.mem_finiteAffine_normal_iff, FalseTateExtension.finiteAffine_units_sup_index, FalseTateExtension.finiteAffine_normal_index}
\leanok
For a normal subgroup $N\leq A_n$, put
$H=\iota_T^{-1}(N)\leq T_0$ and $V=\rho_U(N)\leq U_0$, where
$\iota_T(b)=(b,1)$ and $\rho_U(b,u)=u$.
Then $(b,u)\in N$ if and only if $b\in H$ and $u\in V$.
Thus $N=H\rtimes V$, and
\[
 [A_n:U_0N]=[T_0:H],\qquad
 [A_n:N]=[U_0:V][T_0:H].
\]
\end{lemma}
\begin{proof}
The prime is odd, so $2$ is a unit in $\mathbb Z/p^n\mathbb Z$.
For $(b,u)\in N$, normality and the actual affine multiplication give
\[
 (0,2)(b,u)(0,2)^{-1}(b,u)^{-1}=(b,1)\in N.
\]
Multiplication by $(-b,1)$ then gives $(0,u)\in N$.
If $u\in V$, choose a preimage $(c,u)\in N$ and apply this argument
to obtain $(0,u)\in N$; if also $b\in H$, their product
$(b,1)(0,u)=(b,u)$ belongs to $N$. This proves the membership assertion.
Conjugating $(b,1)\in N$ by $(0,v)$ for arbitrary $v\in U_0$ shows
that $H$ is stable under every unit. Consequently $U_0N=H\rtimes U_0$:
every product has translation in $H$, and
$(b,v)=(0,v)(v^{-1}b,1)$ proves the reverse inclusion.
The genuine coordinate bijections identify the orders of
$N$, $U_0N$ and $A_n$ with $|H||V|$, $|H||U_0|$ and
$|T_0||U_0|$, respectively. Dividing these finite orders proves
both index formulas.
\end{proof}

\begin{lemma}[The affine shape of the actual upper groups]
\label{ft:finite-upper-structure}
\uses{ft:finite-galois-structure,ft:normal-affine-subgroup,ft:finite-rational-cyclotomic-filtration,fw:finite-tower-upper-restriction}
\lean{FalseTateExtension.finiteUpperAffineGroup, FalseTateExtension.finiteUpperAffine_unit_projection, FalseTateExtension.finiteUpperAffine_unit_interval, FalseTateExtension.finiteUpperAffine_unit_tail}
\leanok
For $n\geq1$ and $u>0$, let $N_n(u)$ be the image of the original
group $G_n^u$ under its actual affine isomorphism, and put
$H_n(u)=\iota_T^{-1}(N_n(u))\leq T_0$. Then $N_n(u)$ is normal and
\[
 N_n(u)=H_n(u)\rtimes U_{\lceil u\rceil}.
\]
\end{lemma}
\begin{proof}
An original upper group is an original lower group with its parameter
changed by the finite Herbrand inverse, so it is normal.
Its image under the actual Galois isomorphism is therefore normal.
Restriction to the original cyclotomic subfield is the unit coordinate
of this same isomorphism. The finite upper-restriction theorem identifies
its image with the original cyclotomic upper group, which the checked
cyclotomic filtration identifies as $U_{\lceil u\rceil}$.
This formula includes every right endpoint and triviality when
$\lceil u\rceil\geq n$. The preceding normal-coordinate splitting now
gives exactly the displayed subgroup.
\end{proof}

\begin{lemma}[The intrinsic radical index measures actual translation depth]
\label{ft:finite-upper-radical-index}
\uses{ft:finite-upper-structure,ft:finite-galois-structure,ft:total-ramification,ft:finite-radical-index-profile,fw:finite-galois-overfield-apf}
\lean{FalseTateExtension.finiteUpperAffine_translation_index}
\leanok
For $n\geq1$ and $u>0$,
\[
 a_{F_n/\mathbb Q_p}(u)
 =[G_n:\operatorname{Gal}(L_n(a)/F_n)G_n^u]
 =[T_0:H_n(u)]=p^{m_n(u)},
\]
where $m_n$ is the strict-threshold count of the radical index profile.
\end{lemma}
\begin{proof}
Set $J=\operatorname{Gal}(L_n(a)/F_n)$, the actual pure-unit subgroup.
The intrinsic finite Galois-overfield formula gives
\[
 a_{F_n/\mathbb Q_p}(u)
 =[(G_n)_0:(J\cap(G_n)_0)G_n^u].
\]
Actual total ramification gives $(G_n)_0=G_n$, proving the first equality.
In the actual affine coordinates, $J$ is $U_0$, and the preceding
normal-coordinate index formula identifies $[G_n:JG_n^u]$ with
$[T_0:H_n(u)]$. The already proved pure-radical index profile says
$a_{F_n/\mathbb Q_p}(u)=p^{m_n(u)}$, giving the last equality.
All indices therefore refer to the original fields and their actual
upper groups.
\end{proof}

\begin{lemma}[The Kummer upper jumps]
\label{ft:kummer-upper-jumps}
\uses{ft:finite-radical-index-profile,ft:finite-upper-radical-index}
\lean{FalseTateExtension.finiteUpperAffine_translation}
\leanok
For $n\geq1$ and $u>0$, the actual translation subgroup is
$H_n(u)=T_{m_n(u)}$. In particular, put $k_n(u)=m_n(u)$.
For $1\leq r\leq n$ and $b_r=r+1/(p-1)$,
\[
 k_n(u)<r\quad\Longleftrightarrow\quad 0<u\leq b_r,\qquad
 k_n(u)=\#\{r:1\leq r\leq n,\ b_r<u\}.
\]
Thus the actual translation depth increases by exactly one immediately
after each $b_r$, including the single jump when $n=1$.
\end{lemma}
\begin{proof}
The actual translation group $T_0$ is cyclic of order $p^n$.
The previous lemma computes $[T_0:H_n(u)]=p^{m_n(u)}$.
The subgroup $T_{m_n(u)}$ has that same index, since
$0\leq m_n(u)\leq n$. Subgroups of a finite cyclic group with the same
index are equal, so $H_n(u)=T_{m_n(u)}$.
The thresholds $b_1,\ldots,b_n$ are strictly increasing.
Fewer than $r$ are strictly less than $u$ precisely when $u\leq b_r$,
which proves the equivalence including its right endpoint.
Crossing one threshold adds exactly that threshold to the count,
proving the final assertion.
\end{proof}

\begin{theorem}[Upper ramification groups of the radical layers]
\label{ft:upper-filtration}
\uses{ft:kummer-upper-jumps,ft:finite-upper-structure,ft:total-ramification}
\lean{FalseTateExtension.finiteUpperAffine_nonpos, FalseTateExtension.finiteUpperAffine_initial, FalseTateExtension.finiteUpperAffine_fractional, FalseTateExtension.finiteUpperAffine_integer, FalseTateExtension.finiteUpperAffine_fractional_integer, FalseTateExtension.finiteUpperAffine_integer_fractional, FalseTateExtension.finiteUpperAffine_final_interval, FalseTateExtension.finiteUpperAffine_eq_bot, FalseTateExtension.finiteUpperAffine_integerParameter_interval}
\leanok\source{viviani-2004-radical-extensions:page-0034,viviani-2004-radical-extensions:page-0035}
\dcref{Viviani, Theorem 6.5 and Corollary 6.10}
Let $n\geq2$. At upper breaks the groups of $L_n(a)/\mathbb Q_p$ are
\[
\begin{aligned}
 G_n^0&=C(p^n)\rtimes U_0,\\
 G_n^1&=C(p^n)\rtimes U_1,\\
 G_n^{i+1/(p-1)}&=C(p^{n-i+1})\rtimes U_{i+1},\\
 G_n^{i+1}&=C(p^{n-i})\rtimes U_{i+1},\\
 G_n^{n-1+1/(p-1)}&=C(p^2),\\
 G_n^{n+1/(p-1)}&=C(p).
\end{aligned}
\]
The two indexed rows have $1\leq i\leq n-2$. On each intervening open
left, closed right interval the group equals its value at the right
endpoint. Beyond the last break it is trivial.
\end{theorem}
\begin{proof}
At $u=0$ total ramification gives the whole group. For $u>0$ the affine
shape lemma gives the cyclotomic factor $U_{\lceil u\rceil}$; it drops
immediately after the integers $1,\ldots,n-1$. The Kummer-jump lemma says
that the translation depth is the number of elements of
$\{r+1/(p-1):1\leq r\leq n\}$ strictly less than $u$.
At $u=1$ this number is zero. At $u=i+1/(p-1)$ it is $i-1$; at
$u=i+1$ it is $i$, since $0<1/(p-1)<1$. At the last two fractional
breaks the cyclotomic factor is trivial and the translation depths are
$n-2$ and $n-1$. These give every displayed group. There are no other
drops in either factor, and both factors are trivial after the last break.
\end{proof}

\subsection{The lower indices and the first layer}

\begin{lemma}[Fractional-break lower indices]
\label{ft:fractional-break-index}
\uses{ft:upper-filtration,fw:ramification-filtration}\lean{FalseTateExtension.inverseHerbrandFunction_falseTate_fractional}
\leanok
\source{viviani-2004-radical-extensions:page-0035}
For $n\geq2$ and $t_i=i+1/(p-1)$, the inverse Herbrand function satisfies
\[
 \psi_n(t_i)=\frac{2p^{2i}+p-1}{p+1}
 \qquad(1\leq i\leq n-1).
\]
\end{lemma}
\begin{proof}
On $0<u\leq1$ the upper group has index $p-1$; on
$1<u\leq t_1$ it has index $(p-1)p$.
Consequently $\psi_n(t_1)=p-1+p=2p-1$, which is the asserted value.
For $1\leq i\leq n-2$, the upper indices on
$(t_i,i+1]$ and $(i+1,t_{i+1}]$ are respectively
$(p-1)p^{2i}$ and $(p-1)p^{2i+1}$.
The interval lengths are $(p-2)/(p-1)$ and $1/(p-1)$, so
\[
 \psi_n(t_{i+1})-\psi_n(t_i)
   =(p-2)p^{2i}+p^{2i+1}=2(p-1)p^{2i}.
\]
The proposed expression has this difference and the computed initial
value; induction proves the formula.
\end{proof}

\begin{lemma}[The corrected integer-break index]
\label{ft:integer-break-index}
\uses{ft:upper-filtration,ft:fractional-break-index,fw:ramification-filtration}
\lean{FalseTateExtension.inverseHerbrandFunction_falseTate_integer}
\leanok\source{false-tate-lift:page-0011}
\dcref{Yuan, Remark 5.2}
For $n\geq2$ and $1\leq i\leq n-2$,
\[
 \psi_n(i+1)=\frac{(p-1)(p^{2i+1}+1)}{p+1}.
\]
\end{lemma}
\begin{proof}
At $t_i=i+1/(p-1)$ the preceding lemma gives
$\psi_n(t_i)=(2p^{2i}+p-1)/(p+1)$. On $(t_i,i+1]$ the index is
$(p-1)p^{2i}$. Integration adds $(p-2)p^{2i}$, and therefore
\[
 \psi_n(i+1)=\frac{2p^{2i}+p-1}{p+1}+(p-2)p^{2i}
            =\frac{(p-1)(p^{2i+1}+1)}{p+1}.
\]
\end{proof}

\begin{lemma}[The final lower break]
\label{ft:final-break-index}
\uses{ft:upper-filtration,ft:fractional-break-index,fw:ramification-filtration}
\lean{FalseTateExtension.inverseHerbrandFunction_falseTate_final, FalseTateExtension.inverseHerbrandFunction_falseTate_one}
\leanok\source{viviani-2004-radical-extensions:page-0035}
For $n\geq2$,
\[
 \psi_n\left(n+\frac1{p-1}\right)
  =\frac{p^{2n}+p^{2n-2}+p-1}{p+1},\qquad
 \psi_n(1)=p-1.
\]
\end{lemma}
\begin{proof}
The second formula integrates the constant index $p-1$ on $(0,1]$.
Between $t_{n-1}=n-1+1/(p-1)$ and $t_n=n+1/(p-1)$ the upper group
is $C(p)$ and has index $(p-1)p^{2n-2}$ in the group of order
$(p-1)p^{2n-1}$. This interval has length one. Hence
\[
 \psi_n(t_n)=\frac{2p^{2n-2}+p-1}{p+1}+(p-1)p^{2n-2}
            =\frac{p^{2n}+p^{2n-2}+p-1}{p+1}.
\]
\end{proof}

\begin{theorem}[Corrected finite ramification filtration]
\label{ft:finite-ramification}
\uses{ft:upper-filtration,ft:fractional-break-index,ft:integer-break-index,ft:final-break-index,fw:ramification-filtration}
\lean{FalseTateExtension.lowerRamificationGroupReal_falseTate_nonpos, FalseTateExtension.lowerRamificationGroupReal_falseTate_initial, FalseTateExtension.lowerRamificationGroupReal_falseTate_fractional, FalseTateExtension.lowerRamificationGroupReal_falseTate_integer, FalseTateExtension.lowerRamificationGroupReal_falseTate_final, FalseTateExtension.lowerRamificationGroupReal_falseTate_fractional_integer, FalseTateExtension.lowerRamificationGroupReal_falseTate_integer_fractional, FalseTateExtension.lowerRamificationGroupReal_falseTate_final_interval, FalseTateExtension.lowerRamificationGroupReal_falseTate_eq_bot, FalseTateExtension.inverseHerbrandFunction_falseTate_integerParameter}
\leanok\source{viviani-2004-radical-extensions:page-0035,false-tate-lift:page-0011}
\dcref{Viviani, Corollary 6.10; Yuan, Remark 5.2}
For $n\geq2$, the upper and lower groups of $L_n(a)/\mathbb Q_p$ satisfy
\[
\begin{aligned}
 G_n^0&=C(p^n)\rtimes U_0=(G_n)_0,\\
 G_n^1&=C(p^n)\rtimes U_1=(G_n)_{p-1},\\
 G_n^{i+1/(p-1)}&=C(p^{n-i+1})\rtimes U_{i+1}
       =(G_n)_{(2p^{2i}+p-1)/(p+1)},\\
 G_n^{i+1}&=C(p^{n-i})\rtimes U_{i+1}
       =(G_n)_{(p-1)(p^{2i+1}+1)/(p+1)},\\
 G_n^{n-1+1/(p-1)}&=C(p^2)
       =(G_n)_{(2p^{2n-2}+p-1)/(p+1)},\\
 G_n^{n+1/(p-1)}&=C(p)
       =(G_n)_{(p^{2n}+p^{2n-2}+p-1)/(p+1)}.
\end{aligned}
\]
The middle two rows have $1\leq i\leq n-2$. On each interval between
successive upper breaks, open at the left and closed at the right, the
group equals its value at the right endpoint; the same holds in lower
numbering between the corresponding lower breaks. Beyond the last break
the groups are trivial. In particular the integer-break exponent is
$2i+1$.
\end{theorem}
\begin{proof}
The upper filtration theorem gives the groups and all upper intervals.
The three preceding index lemmas compute $\psi_n$ at every positive
break. The definition $G_n^u=(G_n)_{\psi_n(u)}$ gives all equalities.
Since $\psi_n$ is strictly increasing and continuous, it sends each
interval between consecutive upper breaks onto the corresponding
interval between lower breaks, preserving endpoint inclusion. The
integer-break computation gives the exponent $2i+1$, correcting the
printed $2i$ in Viviani's corollary as in Yuan's remark.
\end{proof}

\begin{theorem}[The first radical layer]
\label{ft:first-layer-filtration}
\uses{ft:finite-upper-structure,ft:kummer-upper-jumps,ft:total-ramification,fw:ramification-filtration}
\lean{FalseTateExtension.finiteUpperAffine_first_layer_filtration, FalseTateExtension.inverseHerbrandFunction_falseTate_first_layer, FalseTateExtension.inverseHerbrandFunction_falseTate_first_break, FalseTateExtension.lowerRamificationGroupReal_falseTate_first_layer_filtration}
\leanok\source{viviani-2004-radical-extensions:page-0033}
\dcref{Viviani, Theorem 6.3 and Lemma 6.4, first layer}
For $n=1$, put $b=p/(p-1)$. The full filtration is
\[
 G_1^u=\begin{cases}
 C(p)\rtimes(\mathbb Z/p\mathbb Z)^\times,&-1\leq u\leq0,\\
 C(p),&0<u\leq b,\\
 \{1\},&u>b,
 \end{cases}
 \qquad
 (G_1)_t=\begin{cases}
 G_1,&-1\leq t\leq0,\\
 C(p),&0<t\leq p,\\
 \{1\},&t>p.
 \end{cases}
\]
Thus the positive upper break is $p/(p-1)$ and the positive lower break
is $p$.
\end{theorem}
\begin{proof}
For $n=1$ every positive cyclotomic upper group is trivial, since
$U_j=\{1\}$ for $j\geq1$. The actual affine upper-group shape is
therefore $G_1^u=T_{k_1(u)}$ for $u>0$.
The Kummer-jump formula says that $k_1(u)=0$ for $0<u\leq b_1$
and $k_1(u)=1$ for $u>b_1$, where
$b_1=1+1/(p-1)=p/(p-1)$. This proves the positive upper filtration.
Total ramification gives the whole group at all indices from $-1$ to $0$.
On $0<u\leq b_1$, the upper group $C(p)$ has index $p-1$ in $G_1$.
Consequently
\[
 \psi_1(u)=(p-1)u\quad(0\leq u\leq b_1),\qquad
 \psi_1(b_1)=p.
\]
The defining relation $G_1^u=(G_1)_{\psi_1(u)}$, with its strictly
increasing continuous parameter change, now gives exactly the stated
lower groups, including both endpoints and triviality after $p$.
This derives the first-layer result of Viviani from the same intrinsic
index route as the higher layers.
\end{proof}

For the remaining statements take $K=\Qp$, $\pi=p$, and $a=p$; put
$L_n=L_n(p)$, and retain the compatible radicals $\pi_n$ chosen above.

\begin{lemma}[Canonical tame units in the p-adic integers]
\label{ft:padic-tame-units}
\source{serre-1979-local-fields:page-0041,false-tate-lift:page-0003}
\dcref{Simple-root Hensel lifting; torsion-unit choice in Yuan, Lemma 1.2}
\lean{PadicInt.existsUnique_teichmullerLift, PadicInt.teichmuller, PadicInt.toZMod_teichmuller, PadicInt.teichmullerUnits, PadicInt.toZMod_teichmullerUnits, PadicInt.teichmullerUnitsEquivRange, PadicInt.natCard_teichmullerUnits_range, PadicInt.isCyclic_teichmullerUnits_range, PadicInt.exists_teichmullerUnit_orderOf}
\leanok
For each $a\in\mathbb F_p$ there is a unique $[a]\in\mathbb Z_p$
with residue $a$ and $[a]^p=[a]$. These lifts give an injective
multiplicative map preserving zero and one. On units its image is a
cyclic subgroup of $\mathbb Z_p^\times$ of order $p-1$, and reduction
is its inverse isomorphism to $\mathbb F_p^\times$. In particular
there is a unit of exact order $p-1$ whose reduction generates
$\mathbb F_p^\times$.
\end{lemma}
\begin{proof}
For $F(X)=X^p-X$, choose any integral lift $b$ of $a$. Its reduction
satisfies $\overline{F}(a)=a^p-a=0$ and
$\overline{F}'(a)=-1\ne0$. Hensel's simple-root theorem gives a unique
root $[a]$ with this residue. The roots $0,1$ show $[0]=0$ and $[1]=1$.
The product $[a][b]$ has residue $ab$ and satisfies
$([a][b])^p=[a][b]$, so uniqueness gives $[ab]=[a][b]$.
Reduction proves injectivity. If $a\ne0$, then
$[a][a^{-1}]=1$, so $[a]$ is a unit; conversely its residue must be
nonzero if it is a unit. Thus restriction to $\mathbb F_p^\times$
is an injective group homomorphism with reduction as inverse on its
image. The multiplicative group of the finite field is cyclic of order
$p-1$, so its image has the same properties. Lifting a generator
therefore gives a unit of exact order $p-1$ and the stated primitive
residue. This constructs the units explicitly from simple residue roots;
no torsion-unit or lifting section is assumed.
\end{proof}

\begin{definition}[The tame generator]
\label{ft:delta}\uses{ft:galois-structure,ft:tau,ft:padic-tame-units,ft:infinite-cyclotomic-character}\lean{FalseTateExtension.tameGenerator, FalseTateExtension.padicCyclotomicCharacter_tameGenerator, FalseTateExtension.orderOf_tameGenerator, FalseTateExtension.tameGenerator_fixes_kummerExtension, FalseTateExtension.tauPadicPower, FalseTateExtension.continuous_tauPadicPower, FalseTateExtension.tauPadicPower_intCast, FalseTateExtension.tameGenerator_conj_tau, FalseTateExtension.orderOf_tameGenerator_residueCharacter, FalseTateExtension.relativeTameGenerator, FalseTateExtension.relativeTameGenerator_restrictScalars, FalseTateExtension.orderOf_relativeTameGenerator}
\leanok
\source{false-tate-lift:page-0003}
\dcref{Yuan, Lemma 1.2}
Choose a generator $\delta$ of the subgroup of order $p-1$ in the complement
$\Gal(\Qp^{\mathrm{FT}}/\Qp^{\mathrm{Kum}})\simeq\Zp^\times$.
Then $\delta\tau\delta^{-1}=\tau^{\chi(\delta)}$.
\end{definition}
\begin{proof}
For the rational base the actual cyclotomic character has image all of
$\mathbb Z_p^\times$. The proved splitting identifies the actual
complement fixing the radical extension with this unit group. Transport
a generator of the preceding subgroup of order $p-1$ through this
isomorphism to obtain $\delta$. In the actual affine coordinates,
$\tau=(1,1)$ and $\delta=(0,\chi(\delta))$. The multiplication law
$(b,u)(c,v)=(b+uc,uv)$ gives
\[
 (0,\chi(\delta))(1,1)(0,\chi(\delta)^{-1})
 =(\chi(\delta),1).
\]
Here the continuous power $\tau^c$ for $c\in\mathbb Z_p$ is the
actual Kummer element corresponding to $c$ under its proved additive
isomorphism; its integer powers agree with this definition by the
generator specification and continuity. The last displayed pair is
therefore $\tau^{\chi(\delta)}$, proving the relation.
\end{proof}


\begin{lemma}[The finite order of the prescribed Kummer generator]
\label{ft:finite-kummer-generator-order}
\uses{ft:finite-ramification,ft:total-ramification,ft:tau,fw:totally-ramified-uniformizer-generation,fw:generator-evaluation,fw:integer-lower-groups,fw:integral-order-function}
\lean{FalseTateExtension.kummerTauFinite, FalseTateExtension.finiteDiagonalGaloisEquiv_kummerTauFinite, FalseTateExtension.ramificationOrder_kummerTauFinite, FalseTateExtension.addVal_kummerTauFinite_uniformizer, FalseTateExtension.normalizedValuation_kummerTauFinite_displacement, FalseTateExtension.normalizedValuation_kummerTauFinite_div_sub_one, FalseTateExtension.kummerRamificationBreak_zero, FalseTateExtension.kummerRamificationBreak_mul, FalseTateExtension.kummerRamificationBreak_recurrence, FalseTateExtension.kummerRamificationBreak_growth, FalseTateExtension.kummerRamificationBreak_succ_dvd, FalseTateExtension.kummerRamificationBreak_tendsto_real, FalseTateExtension.kummerRamificationBreak_tendsto_padic}
\leanok
\source{viviani-2004-radical-extensions:page-0035,false-tate-lift:page-0011,serre-1979-local-fields:page-0068,serre-1979-local-fields:page-0069}
Let $p$ be odd, let $\pi$ be a uniformizer of $\mathbb Q_p$, and use
its prescribed compatible roots and Kummer generator $\tau$. Put
$D_N=\mathbb Q_p(\zeta_{p^N},\pi^{1/p^N})$, and let $g_{N,k}$ be the
ordinary restriction of $\tau^{p^k}$ to this original field. Define
\[
 d_k=1+2(p-1)\sum_{j=0}^{k}p^{2j}\in\mathbb N.
\]
For every $k\geq0$ and every $N\geq k+2$, its actual normalized
integral ramification order is $d_k+1$, its shifted lower break is $d_k$,
and every uniformizer $\varpi$ of the original integer ring satisfies
\[
 v_{D_N}\bigl(g_{N,k}(\varpi)-\varpi\bigr)=d_k+1,
 \qquad
 v_{D_N}\bigl(g_{N,k}(\varpi)/\varpi-1\bigr)=d_k.
\]
The sequence has the exact properties
\[
 (p+1)d_k=2p^{2k+2}+p-1,\qquad d_0=2p-1,\qquad p\mid d_k+1,
\]
\[
 d_{k+1}=p^2d_k-(p-1)^2\geq(p+1)d_k,\qquad
 d_k\longrightarrow+\infty\text{ in }\mathbb R,
 \qquad
 d_k\longrightarrow\frac{p-1}{p+1}\text{ in }\mathbb Z_p.
\]
The indexing condition is the diagonal exponent $N\geq k+2$.
The smaller level $N=k+1$ belongs to the distinct final row of the
finite table and is not asserted to have this stable break.
\end{lemma}
\begin{proof}
Write $S_k=\sum_{j=0}^{k}p^{2j}$. The telescoping identity
$(p^2-1)S_k=p^{2k+2}-1$ gives
\[
 (p+1)\bigl(1+2(p-1)S_k\bigr)=2p^{2k+2}+p-1.
\]
Thus the displayed rational expression for $d_k$ is an integer, and
the defining sum gives $d_0=2p-1$. Adding its last summand gives
$d_{k+1}=d_k+2(p-1)p^{2k+2}$; eliminating the power with the preceding
identity gives $d_{k+1}=p^2d_k-(p-1)^2$.

The actual Kummer coordinate is additive on the cyclotomic-character
kernel, and $\tau$ has coordinate one. The restriction formulas for
the actual finite Kummer and cyclotomic coordinates consequently give
\[
 e_N(g_{N,k})=(p^k,1)
 \quad\text{in }(\mathbb Z/p^N\mathbb Z)\rtimes
                  (\mathbb Z/p^N\mathbb Z)^\times.
\]
These are ordinary restrictions under the original field inclusions.
Since $k<N$, the translation $p^k$ belongs to $T_k$ but not to
$T_{k+1}$. For the latter assertion reduce modulo $p^{k+1}$:
every element of $T_{k+1}$ becomes zero, whereas $p^k$ does not.

Set $i=k+1$, so $1\leq i\leq N-1$. At
$b_i=i+1/(p-1)$ the actual upper group has translation part $T_{i-1}$,
and $g_{N,k}$ therefore belongs to it. The finite lower table places
this group at
\[
 \psi_N(b_i)=\frac{2p^{2i}+p-1}{p+1}=d_k.
\]
On $(b_i,i+1]$ the translation part is $T_i$, which excludes
$g_{N,k}$. Its full upper index is $(p-1)p^{2i}$, so
\[
 \psi_N(i+1)-\psi_N(b_i)=(p-2)p^{2i}\geq1.
\]
This calculation also applies to $i=N-1$: the unit factor is already
trivial and this interval is contained in the final fractional interval;
the integer $N$ need not be a new break. Hence
$g_{N,k}\in(G_N)_{d_k}$ but $g_{N,k}\notin(G_N)_{d_k+1}$.
The integer lower-group criterion gives
$d_k+1\leq j_{\mathcal O_{D_N}}(g_{N,k})<d_k+2$, and the order is
integer-valued, so it equals $d_k+1$.

The preceding odd-prime total-ramification theorem applies to the actual
field $D_N$. Every one of its uniformizers generates its original
integer ring over $\mathbb Z_p$, by total-ramification uniformizer
generation. Generator evaluation therefore identifies the same
intrinsic order with $v_{D_N}(g_{N,k}(\varpi)-\varpi)$ for every
$\varpi$. Dividing the displacement by $\varpi$, whose valuation is
one, gives the asserted shifted value. This argument contains no
field-of-norms action or ramification-transfer premise.

Reducing $(p+1)d_k=2p^{2k+2}+p-1$ modulo $p$ gives
$d_k\equiv-1\pmod p$. The defining sum gives $d_k\geq p-1>0$.
Since $p^2-p-1\geq p-1$, one has
$(p^2-p-1)d_k\geq(p-1)^2$, and the recurrence implies
$d_{k+1}\geq(p+1)d_k$. In particular $d_k\geq(p+1)^kd_0$, which
proves divergence in $\mathbb R$.
Finally $p+1$ has residue one and is a unit of $\mathbb Z_p$. In this
same p-adic ring the exact closed formula gives
\[
 d_k-\frac{p-1}{p+1}=\frac{2p^{2k+2}}{p+1}.
\]
The right side tends to zero because $|p|_p=p^{-1}<1$. This proves
the stated p-adic limit in the original p-adic topology.
\end{proof}

\begin{proposition}[The actual norm-field displacement of Kummer powers]
\label{ft:norm-field-kummer-displacement}
\uses{ft:finite-kummer-generator-order,ft:strict-apf,ft:tau,fw:norm-field-isomorphism,fw:cofinal-displacement-valuation,fw:norm-limit-uniformizer}
\lean{FalseTateExtension.normLimitValuation_tau_pow_displacement, FalseTateExtension.normLimitValuation_tau_pow_div_sub_one, FalseTateExtension.normLimitTau_infiniteOrder}
\leanok
\source{false-tate-lift:page-0011,viviani-2004-radical-extensions:page-0035}
Let $p$ be odd, let $\pi$ be a uniformizer of $\mathbb Q_p$, and let
$L=\mathbb Q_p^{\mathrm{FT}}$ with its prescribed Kummer generator $\tau$.
Write $X=X_{\mathbb Q_p}(L)$ for the original field of norm-compatible
families, and $X(g)$ for the action obtained by transporting original finite
coordinates along $g$. For every $u\in X$ with normalized additive valuation
one and every $k\geq0$, one has
\[
 v_X\bigl(X(\tau^{p^k})(u)-u\bigr)=d_k+1,\qquad
 v_X\bigl(X(\tau^{p^k})(u)/u-1\bigr)=d_k,
 \qquad d_k=1+2(p-1)\sum_{j=0}^k p^{2j}.
\]
In particular $X(\tau)$ has infinite order as an automorphism of the same
original norm field.
\end{proposition}
\begin{proof}
The original diagonal fields $D_N=\mathbb Q_p(\zeta_{p^N},\pi^{1/p^N})$
are cofinal among finite subfields of $L$: choose finitely many field generators
of a given finite subfield, put each in a diagonal field, and take the maximum
of their finitely many indices. Their ordinary inclusions identify their
images in $L$ with $D_N$ itself. Each $D_N$ is Galois, so every automorphism
of $L/\mathbb Q_p$ preserves this image. The defining coordinate formula for
$X(g)$ therefore gives
\[
 (X(g)u)_{D_N}=g|_{D_N}(u_{D_N}).
\]
The stable normalized norm-field valuation equals its coordinate valuation on
an actual finite tail. Consequently $u_{D_N}$ has normalized valuation one
for every sufficiently large $N$. It lies in the original spectral integer
ring and is a uniformizer there. For $N\geq k+2$, the finite Kummer-generator
order theorem thus gives
\[
 v_{D_N}\bigl((X(\tau^{p^k})u)_{D_N}-u_{D_N}\bigr)=d_k+1.
\]
To pass from this ordinary coordinate difference to the actual norm-field
difference, fix an integer $s>d_k+1$. Approximate additivity on an actual finite
tail says that their difference belongs to $\mathfrak m_{D_N}^{s}$. Choose
$N$ above that tail, above the stable valuation tail, and above $k+2$.
Its error has valuation strictly greater than $d_k+1$, so the ultrametric
inequality gives exact equality of the valuations of the two differences.
Stable coordinate evaluation now yields the first displayed equality on the
original $X$. Since $v_X(u)=1$, division by $u$ gives the second equality.
Finally the sums defining $d_k$ are strictly increasing because $p>1$.
If two powers $X(\tau)^{p^j}$ and $X(\tau)^{p^k}$ were equal, evaluating at
any norm-field uniformizer and using their exact displacement valuations
would give $d_j=d_k$, hence $j=k$. There are infinitely many distinct powers
of $X(\tau)$, whereas an automorphism of finite order has only finitely many
powers. This proves infinite order.
\end{proof}

\begin{theorem}[The actual false Tate action in Laurent coordinates]
\label{ft:false-tate-laurent-coordinates}
\source{false-tate-lift:page-0002}
\dcref{Yuan, Definition 1.1, the evaluation coordinate}
\uses{ft:strict-apf,ft:total-ramification,ft:infinite-cyclotomic-character,fw:abstract-extension-residue-field,fw:source-residue-constants-field-embedding,fw:uniformizer-laurent-coordinates,fw:laurent-power-series-substitution,fw:norm-field-isomorphism}
\lean{FalseTateExtension.falseTateResidueEquiv, FalseTateExtension.falseTateNormFieldLaurentEquiv, FalseTateExtension.falseTateNormFieldLaurentEquiv_unique, FalseTateExtension.falseTateLaurentAction, FalseTateExtension.falseTateLaurentAction_powerSeries, FalseTateExtension.falseTatePowerSeriesCoordinate_constantCoeff, FalseTateExtension.falseTatePowerSeriesCoordinate_linear_isUnit, FalseTateExtension.falseTatePowerSeriesCoordinate_one, FalseTateExtension.falseTatePowerSeriesCoordinate_mul, FalseTateExtension.falseTatePowerSeriesCoordinate_pow, FalseTateExtension.valuation_falseTatePowerSeriesCoordinate_displacement}
\leanok
Let $p$ be odd, let $\pi$ be a uniformizer of $\mathbb Q_p$, and put
$L=\mathbb Q_p(\mu_{p^\infty},\pi^{1/p^\infty})$ and $X=X_{\mathbb Q_p}(L)$,
where $X$ is the original field of norm-compatible families.
The canonical residue field $k_L$ of the integral closure in $L$ is
canonically isomorphic to $\mathbb F_p$. For every original $u\in X$
with $w(u)=\exp(-1)$ there is a unique continuous field isomorphism
\[
 e_u:\mathbb F_p((T))\xrightarrow{\ \sim\ }X
\]
sending $T$ to $u$ and sending each constant to its actual source
coefficient in $X$. It and its inverse preserve the original normalized
valuations and valuation uniformities.

For $g\in\operatorname{Aut}_{\mathbb Q_p}(L)$, let $X(g)$ be the
original finite-coordinate action and put
$\sigma_{g,u}=e_u^{-1}X(g)e_u$. This is a continuous
$\mathbb F_p$-linear field automorphism. There is a unique
$f_{g,u}\in T\mathbb F_p[[T]]$ with nonzero linear coefficient such that
\[
 \sigma_{g,u}(T)=f_{g,u},\qquad
 \sigma_{g,u}(A)=A\circ f_{g,u}\quad(A\in\mathbb F_p[[T]]).
\]
These are the coordinates of the actual action:
\[
 f_{gh,u}=f_{h,u}\circ f_{g,u},\qquad f_{1,u}=T.
\]
In particular $f_{g^n,u}$ is the $n$-fold compositional iterate of
$f_{g,u}$ for every $n\in\mathbb N$.
\end{theorem}
\begin{proof}
Every finite subextension of $L/\mathbb Q_p$ is totally ramified.
Every residue class of the original infinite algebraic extension has a
representative in an actual finite subextension. The canonical map from
the residue field of $\mathbb Q_p$ to the spectral residue field of $L$
is consequently surjective, and it is injective because it is a unital
map of fields. The identity-valued isomorphism from the canonical integral
closure to the spectral integer ring identifies their residue fields.
Composing these maps with the usual residue isomorphism
$\mathbb Z_p/(p)\simeq\mathbb F_p$ gives the asserted canonical
identification of $k_L$.

The extension $L/\mathbb Q_p$ is infinite: it contains cyclotomic
subextensions of unbounded finite degree. It is APF by strict APF.
The actual source coefficient embedding and its induced residue
isomorphism, transported along $k_L\simeq\mathbb F_p$, therefore satisfy
the hypotheses of the Laurent-coordinate theorem for the original
complete norm field and the specified original uniformizer $u$.
That theorem constructs $e_u$, gives its exact valuation and both
uniform continuities, and proves uniqueness from constants and $T$.

The original norm-field action preserves its actual normalized valuation
and is a group homomorphism. Conjugation by $e_u$ thus gives continuous
valuation-preserving automorphisms satisfying
$\sigma_{gh,u}=\sigma_{g,u}\sigma_{h,u}$. Each fixes $\mathbb F_p$:
a unital ring homomorphism fixes the prime field, since it fixes $1$ and
therefore all integral multiples of $1$.
The Laurent element $\sigma_{g,u}(T)$ has normalized additive valuation
one. Its coefficients in degrees below one vanish, and its coefficient
in degree one is nonzero. It is therefore the image of a unique power
series $f_{g,u}$ with the asserted constant and linear coefficients.

On polynomials in $T$, the equality
$\sigma_{g,u}(A)=A\circ f_{g,u}$ follows by the ring laws, the constant
formula and the value at $T$. Polynomial truncations converge to each
power series. Substitution by a power series with zero constant term is
continuous for these topologies, as follows degree by degree because the
coefficient in degree $m$ depends only on the first $m+1$ input
coefficients. Continuity of both sides extends the equality to every
power series. Applying the group identity to $T$ now gives
$f_{gh,u}=\sigma_{g,u}(f_{h,u})=f_{h,u}\circ f_{g,u}$.
The identity gives $f_{1,u}=T$, and induction gives the iterate formula.
\end{proof}

\begin{lemma}[The actual action profile in every norm-field coordinate]
\label{ft:norm-field-action-profile}
\uses{ft:norm-field-kummer-displacement,ft:false-tate-laurent-coordinates,ft:finite-kummer-generator-order}
\lean{FalseTateExtension.falseTatePowerSeriesCoordinate_tau_pow_valuation, FalseTateExtension.falseTatePowerSeriesCoordinate_tau_pow_order, FalseTateExtension.falseTatePowerSeriesCoordinate_tau_coeff_one, FalseTateExtension.falseTatePowerSeriesCoordinate_tau_coeff_first_ne_zero, FalseTateExtension.falseTatePowerSeriesCoordinate_tau_infiniteOrder, FalseTateExtension.falseTatePowerSeriesCoordinate_tau_iterate_order}
\leanok
\source{false-tate-lift:page-0011,false-tate-lift:page-0012}
Let $p$ be odd, let $\pi$ be a uniformizer of $\mathbb Q_p$, and use
the prescribed original false Tate field, generator $\tau$, and original
norm field $X$. For every original uniformizer $u$ of $X$ put
$f=f_{\tau,u}$. The literal power-series coordinates satisfy
\[
 v_T(f_{\tau^{p^k},u}-T)=d_k+1,\qquad
 d_k=1+2(p-1)\sum_{j=0}^k p^{2j}.
\]
In particular $f=T+\alpha T^{2p}+O(T^{2p+1})$ with
$\alpha\in\mathbb F_p^\times$, and $f$ has infinite compositional order.
All coordinates here are formed from the same original $u$ and the same
original coefficient section.
\end{lemma}
\begin{proof}
The original Laurent evaluation isomorphism $e_u$ preserves normalized
valuations and sends $T$ to $u$. It sends the actual coordinate
$f_{\tau^{p^k},u}$ to $X(\tau^{p^k})(u)$. Its additivity and the actual
norm-field Kummer displacement theorem therefore give
\[
 \exp\bigl(-v_T(f_{\tau^{p^k},u}-T)\bigr)
 =w\bigl(X(\tau^{p^k})(u)-u\bigr)=\exp(-(d_k+1)).
\]
The exponential is injective and its value is nonzero, so the displacement
is nonzero and has the stated exact integer order. The first coefficient
at that order is consequently nonzero and every smaller displacement
coefficient vanishes. At $k=0$, $d_0+1=2p\geq2$, giving the displayed
expansion of $f$ with tangent coefficient one.
The actual anti-composition law identifies $f_{\tau^{p^k},u}$ with the
$p^k$-fold iterate of $f$. If $f$ had a positive compositional period $m$,
all its iterates would belong to the finite list
$T,f,\ldots,f^{\circ(m-1)}$, by dividing each exponent by $m$.
Their orders relative to $T$ would then take only finitely many values.
But the orders $d_k+1$ are strictly increasing, since each successive
defining sum adds the positive summand $2(p-1)p^{2k+2}$.
This contradiction proves infinite compositional order.
\end{proof}

\begin{proposition}[Ramification numbers on the field of norms]
\label{ft:ramification-numbers}
\uses{ft:norm-field-action-profile,ft:norm-field-kummer-displacement,ft:finite-kummer-generator-order,ft:finite-ramification,ft:strict-apf,ft:total-ramification,ft:tau,ft:false-tate-laurent-coordinates}
\lean{FalseTateExtension.falseTatePowerSeriesCoordinate_tau_pow_order, FalseTateExtension.kummerRamificationBreak_zero, FalseTateExtension.kummerRamificationBreak_mul, FalseTateExtension.kummerRamificationBreak_recurrence, FalseTateExtension.kummerRamificationBreak_growth, FalseTateExtension.kummerRamificationBreak_succ_dvd, FalseTateExtension.kummerRamificationBreak_tendsto_real, FalseTateExtension.kummerRamificationBreak_tendsto_padic}
\leanok
\source{false-tate-lift:page-0011}
\dcref{Yuan, Proposition 5.1}
Put $E=X_{\Qp}(\Qp^{\mathrm{FT}})$. For every uniformizer $u$ of $E$,
write $f_{g,u}=\operatorname{ev}_u^{-1}(g(u))$ and set
$d_k=v_T(f_{\tau^{p^k},u}-T)-1$. Then
\[
 d_k=\frac{2p^{2k+2}+p-1}{p+1},\qquad
 d_0=2p-1,\qquad d_{k+1}=p^2d_k-(p-1)^2\geq(p+1)d_k.
\]
In particular $d_k\equiv-1\pmod p$, $d_k\longrightarrow+\infty$ in
$\mathbb R$, and $d_k\longrightarrow(p-1)/(p+1)$ in $\Zp$.
\end{proposition}
\begin{proof}
The preceding actual profile identifies the shifted order with the same
integer $d_k$ computed in the finite Kummer-generator theorem.
That theorem gives $(p+1)d_k=2p^{2k+2}+p-1$ and all the displayed
recurrence, growth, divisibility and real and p-adic limit identities.
Since $p+1\ne0$, dividing the closed identity gives the asserted formula.
\end{proof}

\begin{lemma}[The tame action in any coordinate]
\label{ft:tame-action}\uses{fw:norm-field-action-continuity,ft:ramification-numbers,ft:delta}
\lean{FalseTateExtension.continuous_falseTatePowerSeriesCoordinate_tauPadicPower_coeff, FalseTateExtension.falseTatePowerSeriesCoordinate_tauPadicPower_coeff_displacement, FalseTateExtension.falseTatePowerSeriesCoordinate_tameGenerator_coeff_one, FalseTateExtension.orderOf_falseTatePowerSeriesCoordinate_tameGenerator_linear, FalseTateExtension.orderOf_falseTatePowerSeriesCoordinate_tameGenerator_linear_unit}
\leanok
\source{false-tate-lift:page-0011,false-tate-lift:page-0012}
\dcref{Yuan, Lemma 5.3}
For every uniformizer $u$ of $E$,
\[
 f_{\delta,u}=\overline{\chi(\delta)}T+O(T^2).
\]
Its linear coefficient has order $p-1$ in $\Fp^\times$.
\end{lemma}
\begin{proof}
Fix the original uniformizer $u$, and write $f=f_{\tau,u}$,
$w=f_{\delta,u}$, $\lambda=w'(0)\in\mathbb F_p^\times$ and
$N=d_0+1=2p$. The actual action profile gives
$f=T+\alpha T^N+O(T^{N+1})$ with $\alpha\ne0$.
We first record the finite coefficient calculation that will be used.
If $g=T+\beta T^N+O(T^{N+1})$, then for every $i\geq1$,
\[
 g^i-T^i=(g-T)\sum_{j=0}^{i-1}g^jT^{i-1-j}
       =O(T^{N+i-1}).
\]
Consequently $g^N=T^N+O(T^{N+1})$ because $N\geq2$, and substitution
in $f$ gives
$f\circ g=T+(\alpha+\beta)T^N+O(T^{N+1})$.
Induction, including the identity iterate, gives
\[
 f^{\circ n}=T+\overline n\alpha T^N+O(T^{N+1})
 \quad(n\in\mathbb N).
\]

For $a\in\mathbb Z_p$ let $h_a=f_{\tau^a,u}$, using the actual continuous
Kummer powers in the original Galois group. The original orbit map at
the same $u$ is continuous by the finite-fixer continuity theorem.
Composing with the original continuous inverse Laurent coordinate map
and each continuous Laurent coefficient map proves continuity of
$a\mapsto[T^j]h_a$ for every $j$. For $a=n\in\mathbb N$, the actual
power and anti-composition laws identify $h_n$ with $f^{\circ n}$.
Natural numbers are dense in $\mathbb Z_p$. Equality of the continuous
coefficient maps to the Hausdorff discrete field $\mathbb F_p$ therefore
extends the preceding formula, coefficient by coefficient through degree
$N$, to every $a$:
\[
 h_a=T+\overline a\alpha T^N+O(T^{N+1}).
\]

Take $a=\chi(\delta)$. The original Galois relation
$\delta\tau\delta^{-1}=\tau^a$, equivalently
$\delta\tau=\tau^a\delta$, and the actual anti-composition law give
$w\circ h_a=f\circ w$.
To compare their first changed coefficients, write
$w=\lambda T+\sum_{i\geq2}c_iT^i$.
The displayed factorization with $g=h_a$ shows that every
$h_a^i-T^i$ with $i\geq2$ has order at least $N+1$. Hence
\[
 [T^N](w\circ h_a-w)=\lambda\overline a\alpha.
\]
On the other side,
$f\circ w-w=\alpha w^N+O(T^{N+1})$, and
$w^N=\lambda^NT^N+O(T^{N+1})$, so
\[
 [T^N](f\circ w-w)=\alpha\lambda^N.
\]
Cancelling the nonzero factors $\alpha$ and $\lambda$ yields
$\lambda^{N-1}=\overline a$.
Since $\lambda^{p-1}=1$ and $N-1=2(p-1)+1$, the left side is
$\lambda$. Thus $\lambda=\overline{\chi(\delta)}$, as asserted.
The actual prescribed tame generator has cyclotomic residue character
of exact order $p-1$, so this same equality gives the asserted order of
the linear coefficient.
\end{proof}


These results use the Galois structure from \cite{Liu2008}, the finite
ramification calculation of \cite{Viviani2004}, and its correction in
\cite{Yuan}. The strict APF criterion is taken from \cite{Wintenberger1983}
and \cite{Sen1972}. The coefficient coordinate is always formed using the
same uniformizer as the evaluation map.
